% concatenated from neurips_2024.tex
 \documentclass{article}


% if you need to pass options to natbib, use, e.g.:
%     \PassOptionsToPackage{numbers, compress}{natbib}
% before loading neurips_2024


% ready for submission
% \usepackage{neurips_2025}


% to compile a preprint version, e.g., for submission to arXiv, add add the
% [preprint] option:
    \usepackage[final]{neurips_2025}


% to compile a camera-ready version, add the [final] option, e.g.:
%     \usepackage[final]{neurips_2024}


% to avoid loading the natbib package, add option nonatbib:
%    \usepackage[nonatbib]{neurips_2024}


\usepackage[utf8]{inputenc} % allow utf-8 input
\usepackage[T1]{fontenc}    % use 8-bit T1 fonts
\usepackage{hyperref}       % hyperlinks
\usepackage{url}            % simple URL typesetting
\usepackage{booktabs}       % professional-quality tables
\usepackage{amsfonts}       % blackboard math symbols
\usepackage{nicefrac}       % compact symbols for 1/2, etc.
\usepackage{microtype}      % microtypography
\usepackage{xcolor}         % colors
\usepackage{enumitem}
\usepackage{arydshln}
\usepackage{graphicx} 
\usepackage{float}
\usepackage{algorithm, algorithmic}
\usepackage{pifont}
% \newtheorem{definition}{Definition}
\newcommand{\cmark}{\textcolor{green!60!black}{\ding{51}}} % checkmark
\newcommand{\xmark}{\textcolor{red!70!black}{\ding{55}}}   % cross


\usepackage{amsthm}

% Define theorem-like environments
\newtheorem{definition}{Definition}
\theoremstyle{remark}
\newtheorem{remark}{Remark}


\title{Causal LLM Routing: End-to-End Regret Minimization from Observational Data}


% The \author macro works with any number of authors. There are two commands
% used to separate the names and addresses of multiple authors: \And and \AND.
%
% Using \And between authors leaves it to LaTeX to determine where to break the
% lines. Using \AND forces a line break at that point. So, if LaTeX puts 3 of 4
% authors names on the first line, and the last on the second line, try using
% \AND instead of \And before the third author name.


\author{%
  Asterios Tsiourvas\\
  % Operations Research Center\\
  MIT % \\
  % Cambridge, MA 02142 \\
  % \texttt{atsiour@mit.edu} \\
  % examples of more authors
  \And
  Wei Sun\\
   IBM Research %\\
  % Yorktown Heights, NY \\
  % \texttt{sunw@us.ibm.com} \\
  \And
  Georgia Perakis \\
  % Operations Research Center\\
  MIT
  % \\
  % Cambridge, MA 02142 \\
  % \texttt{georgiap@mit.edu} \\
  % \And
  % Coauthor \\
  % Affiliation \\
  % Address \\
  % \texttt{email} \\
  % \And
  % Coauthor \\
  % Affiliation \\
  % Address \\
  % \texttt{email} \\
}

\usepackage{amsmath} 
\usepackage{amsthm}
\usepackage{bbm}
\newtheorem{assumption}{Assumption}
\newtheorem{proposition}{Proposition}
\newtheorem{corollary}{Corollary}

\begin{document}


\maketitle


\begin{abstract}
LLM routing aims to select the most appropriate model for each query, balancing competing performance metrics such as accuracy and cost across a pool of language models. Prior approaches typically adopt a \emph{decoupled} strategy, where metrics are first predicted and the model is then selected based on these estimates. This setup is prone to compounding errors and often relies on \emph{full-feedback} data, where each query is evaluated by all candidate models, which is costly to obtain and maintain in practice. In contrast, we learn from \emph{observational data}, which records only the outcome of the model actually deployed. We propose a \emph{causal end-to-end} framework that learns routing policies by minimizing decision-making regret from observational data. To enable efficient optimization, we introduce two theoretically grounded surrogate objectives: a classification-based upper bound, and a softmax-weighted regret approximation shown to recover the optimal policy at convergence. We further extend our framework to handle heterogeneous cost preferences via an interval-conditioned architecture. Experiments on public benchmarks show that our method outperforms existing baselines, achieving state-of-the-art performance across different embedding models. 
\end{abstract}

\section{Introduction}

%LLM routing is an emerging area of research that aims to optimize the selection and utilization of large language models (LLMs) for diverse tasks. The performance of LLMs is highly dependent on both the task and the specific input query \cite{hu2024routerbench}, making model selection a critical challenge. Routing strategies address this by dynamically selecting the most appropriate model for a given input \cite{shnitzer2023large}. Importantly, LLMs vary significantly in their computational demands and associated costs. For example, the cost per one million output tokens ranges from \$1 for LLaMA-3-70B to \$60 for GPT-4-0613, \$1.25 for Haiku, and \$75 for Opus \cite{ong2406routellm}. Generally, lower-cost models are less computationally intensive, which can lead to faster inference times and improved user experience, albeit sometimes at the expense of output quality. Effective routing seeks to balance performance and cost, enabling organizations to reduce operational expenses while maintaining satisfactory task performance.  

LLM routing is an emerging research area focused on optimizing model selection for each input query, balancing performance and cost across a pool of available LLMs. Because LLM performance varies significantly by task and input~\citep{hu2024routerbench}, as well as computational cost~\citep{ong2406routellm}, dynamic routing strategies have been proposed to select the most suitable model per query~\citep{shnitzer2023large}. This challenge becomes even more critical in agentic applications, where multiple LLM calls may be made within a single workflow, making efficient model selection essential for user experience and resource allocation. As LLM deployment scales, routing also contributes to environmental sustainability by reducing unnecessary computation~\citep{singh2025survey}.

%There are two main paradigms of LLM routing. The first is \emph{non-predictive routing} \cite{wang2023fusing,chen2023frugalgpt}, where the router repeatedly calls and evaluates the utilitys of all available LLMs to pick the best for a given query and task. The disadvantage of this paradigm is that non-predictive routing requires evaluating multiple LLMS for all queries, which is typically time-consuming and expensive. As an alternative, researchers have been actively considering \emph{predictive routing} \cite{shnitzer2023large}, in which LLM queries are used as input to a machine learning model that predicts the most suitable LLM for the specific query.\\

%There are two main paradigms in LLM routing. In \emph{non-predictive routing} \citep{wang2023fusing, chen2023frugalgpt}, LLMs are queried sequentially in a cascade, i.e., starting with cheaper, less capable models and escalating to more powerful ones as needed. While straightforward, this approach requires querying multiple models per input, leading to high latency and cost. In contrast, \emph{predictive routing} \citep{shnitzer2023large} uses a learned model to predict the most suitable LLM for a given query upfront, enabling faster and more cost-efficient decisions.


Routing methods can be broadly classified based on whether they invoke one or multiple models per query. \emph{Multi-model} approaches include \emph{non-predictive routing}, which cascades models sequentially from light to heavy~\citep{wang2023fusing, chen2023frugalgpt}, and \emph{predictive ensembles}, which learn to combine outputs or scores from multiple LLMs~\citep{shekhar2024towards, huang2025thriftllm}. While these methods can improve accuracy and robustness, they incur significant computational cost and latency due to repeated inference.
By contrast, \emph{predictive routing} methods, including our approach, select a single LLM for each query by training a router that maps input queries to the most appropriate model~\citep{shnitzer2023large, ong2406routellm, somerstep2025carrot}. This framework offers a more scalable and cost-efficient solution, particularly in settings where minimizing latency and compute resources is critical.


%A common formulation in predictive routing is to maximize a utility function of the form \( y = \text{accuracy} - \lambda \times \text{cost} \), where \( \lambda \geq 0 \) denotes the user's cost sensitivity or willingness-to-pay. Moreover, existing work often follows a \emph{decoupled approach}, where separate predictors are trained to estimate the expected accuracy and cost for each candidate model, conditioned on the input query. A routing decision is then made by evaluating the utility for each model and selecting the one with the highest predicted value. However, the quality of the decision is highly sensitive to the accuracy of both predictors, and errors can compound, leading to suboptimal outcomes. This challenge becomes even more pronounced when additional metrics (e.g., latency) are incorporated, further increasing modeling complexity and amplifying prediction uncertainty.


%A common formulation in predictive routing is to maximize a utility function of the form $y_x(t) = a_x(t) - \lambda \cdot c_x(t)$, where $a_x(t)$ and $c_x(t)$ denote the predicted accuracy and cost of a candidate model $t$ for query $x$, and $\lambda \geq 0$ represents the user's cost sensitivity or willingness-to-pay. Existing work typically follows a \emph{decoupled} approach, where separate predictors are trained for $a_x(t)$ and $c_x(t)$, conditioned on the input query. A routing decision is then made by computing the utility for each model and selecting the one with the highest predicted value. However, the decision quality is highly sensitive to both predictors, and compounding errors can lead to suboptimal routing. This challenge becomes even more pronounced when additional metrics (e.g., latency, faithfulness, alignment) are incorporated, further increasing modeling complexity and amplifying prediction uncertainty. 

A common formulation in predictive routing is to maximize a utility function of the form \( y_x(t) = a_x(t) - \lambda \cdot c_x(t) \), where \( a_x(t) \) and \( c_x(t) \) denote the quality (e.g., accuracy) and cost for model \( t \) on a given query $x$, and \( \lambda \geq 0 \) captures user cost sensitivity, or willingness-to-pay. Existing methods typically adopt a \emph{decoupled} approach: separate predictors are trained for each metric, and routing is performed by selecting the model with the highest estimated utility. However, decision quality is highly sensitive to these predictors, and errors can compound, especially when incorporating additional metrics (e.g., latency, faithfulness, alignment), increasing complexity and uncertainty.

%Another limitation of existing predictive routing methods is their reliance on \emph{full-feedback} datasets, where each query is evaluated by all available LLMs~\citep{ong2406routellm, somerstep2025carrot}. This assumption is impractical due to the high cost of exhaustive evaluation and the rapid evolution of LLMs. In contrast, \emph{observational data}, where each query is routed to a single model and labeled with the resulting output quality, is more scalable and readily available in real-world deployments. However, it introduces \emph{treatment bias} from historical routing policies, which can lead to suboptimal decisions if not properly addressed~\citep{swaminathan2015batch, kunzel2019metalearners}.


%Another fundamental limitation in the existing literature is the reliance on a hyperparameter (typically $\lambda$), also known as a budget constraint, to balance performance and cost \cite{somerstep2025carrot}. Specifically, most works formulate the predictive routing problem as maximizing the weighted objective function $y=accuracy - \lambda \times cost$, where $\lambda \geq 0$ is a hyperparameter with units of accuracy per dollar. Such methods typically model accuracy and cost independently, which can result in inconsistencies and poor decisions. This challenge is exacerbated in more realistic scenarios where additional constraints, such as latency, are introduced—for instance, \(y = \text{accuracy} - \lambda \times \text{cost} - \mu \times \text{latency}\), with \(\mu \geq 0\). In these cases, prior methods may need to predict three separate quantities (accuracy, cost, latency), increasing modeling complexity and compounding prediction errors, ultimately degrading decision quality.


Another fundamental limitation in the existing literature on predictive routing is its reliance on \emph{full-feedback} datasets. Prior work \citep{ong2406routellm} \citep{somerstep2025carrot} assumes access to data where each query has been evaluated by all available LLMs. This assumption is impractical: (i) the computational and monetary cost of exhaustively querying all models is prohibitive; and (ii) the rapid pace of LLM development makes it challenging to maintain comprehensive, up-to-date evaluation datasets. In contrast, \emph{observational data}, where each query is evaluated by only one model, is readily available from real-world LLM deployments, making it far more scalable than full feedback datasets. However, it introduces \emph{treatment bias} from historical routing policies, which can lead to suboptimal decisions if not properly addressed~\citep{swaminathan2015batch, kunzel2019metalearners}.%However, a key challenge with observational data is \emph{treatment bias} which arises from historical policy governing model selection and can lead to biased or suboptimal routing decisions if not properly addressed~\citep{swaminathan2015batch,kunzel2019metalearners}.




%\paragraph{Motivating Example.} Consider two LLMs with accuracy--cost pairs: A = (0.91, 0.50), B = (0.90, 0.05), and a user-defined trade-off parameter \(\lambda = 1\). The true utilitys are \(y_A = 0.41\) and \(y_B = 0.85\), making LLM B the optimal choice. Suppose we train separate models to predict accuracy and cost. While accuracy is often stable across queries, cost can be noisy and context-dependent (e.g., due to input length or system load). A misestimate—\(\hat{a}_A = 0.89\), \(\hat{c}_A = 0.38\); \(\hat{a}_B = 0.88\), \(\hat{c}_B = 0.40\)—leads to predicted utilitys \(\hat{y}_A = 0.51\), \(\hat{y}_B = 0.48\), and an incorrect routing to A. Despite small prediction errors, especially in accuracy, this results in a high regret of \(0.44\). This highlights a key challenge: even minor inaccuracies in cost prediction—common in practice—can lead to suboptimal decisions. 





%Our work departs from prior methods by leveraging \emph{observational data} to learn an \emph{end-to-end} routing policy, enabling accurate and efficient decision-making in practical deployments. 
To the best of our knowledge, this is the first work to (i) learn LLM routing from {observational data}  and (ii) introduce an integrated learning framework for routing. %that avoids the pitfalls of decoupled approaches. 
Our main contributions are:


\begin{itemize}
\item We propose a \emph{causal end-to-end} framework that learns routing policies by directly minimizing decision-making regret using {observational data}. Unlike the predominant {decoupled} paradigm, where various performance metrics (e.g., accuracy, cost) are first predicted and then used to inform routing decisions, our method integrates prediction and prescription into a unified objective. By optimizing for regret directly, the framework is explicitly aligned with the final routing decision quality. Furthermore, it is designed to scale efficiently and leverage readily available observational data, while accounting for treatment bias without requiring costly full-feedback datasets.

\item As the original regret minimization objective is not directly differentiable, we derive two \emph{surrogate objectives} to enable end-to-end policy learning. The first is a {classification-based upper bound} that reframes regret minimization as a multiclass prediction problem under mild Lipschitz assumptions, allowing efficient training with standard methods. The second is a {softmax-weighted regret surrogate} that smoothly approximates regret using a softmax distribution and provably recovers optimal decisions at convergence.

\item We extend our framework to support \emph{heterogeneous cost preferences} by introducing a unified model that conditions on both the query and the user’s cost sensitivity. Leveraging the affine structure of the utility function, we design an efficient interpolation scheme using only two endpoint models per interval. We theoretically show that the optimal treatment is piecewise constant in the cost parameter and that our architecture can exactly represent the optimal policy, enabling flexible and scalable routing across diverse preferences.

\item We conduct \emph{comprehensive experiments} on two public benchmarks, demonstrating that our regret-minimizing and heterogeneous cost-aware approaches consistently outperform existing baselines. Our methods achieve state-of-the-art performance across both BERT-based and LLaMA-based embeddings, highlighting their robustness and practical effectiveness.



    %\item We introduce a \emph{causal end-to-end} framework that learns routing policies by directly minimizing decision-making regret, enabling causal reasoning from {observational data}. In contrast to the predominant two-step paradigm in causal machine learning—where models first estimate outcomes and then optimize over predicted values—our approach is end-to-end. The standard pipeline often suffers from a disconnect between prediction and decision-making, leading to compounding errors and suboptimal prescriptions. Instead, we propose a unified training objective that directly minimizes regret using the doubly robust estimator from \cite{dudik2011doubly}, ensuring that the model is explicitly optimized for the final decision quality. This end-to-end framework is \emph{broadly applicable beyond LLM routing and can be employed in any causal decision-making problem with partial feedback}. Unlike prior work, our objective is differentiable and can be seamlessly incorporated into standard deep learning pipelines.


    
    %\item To avoid the need to retrain separate models for each trade-off parameter \(\lambda\), we develop a \emph{multi-task, budget-aware router} that generalizes across user-specific preferences. By conditioning the routing model on a continuous budget parameter, our approach enables a single policy to dynamically adapt to varying accuracy--cost trade-offs at inference time. To support this, we introduce a novel \emph{interpolation-based architecture} that reuses pre-trained models at selected budget levels and smoothly generalizes between them. We provide theoretical guarantees showing that interpolation yields exact utility reconstruction thereby offering both computational efficiency and provable correctness.

    %\item We conduct comprehensive evaluations on two real-world benchmarks---RouterBench and SPROUT---demonstrating that our regret-minimizing and budget-aware approaches consistently outperform existing routing baselines. Our methods achieve state-of-the-art performance across both BERT-based and LLaMA-based embeddings, validating their robustness and practical applicability.
\end{itemize}






% \section{Related Work}

\section{Methodology}

\subsection{Problem Formulation}

We consider a dataset of $n$ observational samples, denoted by $ \mathcal{D} = \{(x_i, t_i, a_i, c_i)\}_{i=1}^n $, where each sample is independently drawn from the joint distribution $ p(x, t, a, c) $. Here, $ x_i \in \mathcal{X} \subset \mathbb{R}^d $ denotes a feature vector, typically an embedding, that characterizes the query; 
$ t_i \in [\mathcal{T}] := \{0, 1, \dots, T-1\} $ specifies the LLM assigned to the query; %$a_i \in [0, 1] $ records the accuracy of the model's response; 
$a_i \in \mathbb{R}_{\geq 0}$ denotes a numeric quality score of the LLM's response, such as accuracy, or a preference rating; and $ c_i \in \mathbb{R}_{\geq 0} $ represents the cost incurred by model $t_i $ when processing query $x_i$. 



Given a query $x \in \mathcal{X}$, the objective of an LLM router is to select a model $t \in \mathcal{T}$ that maximizes the cost-aware performance, or utility, defined as $y_x(t) := a_x(t) - \lambda \cdot c_x(t)$. Here, $\lambda \geq 0$ is a user-specified parameter, modeling the trade-off between accuracy and cost. Higher values of $y_x(t)$, corresponding to greater performance, are preferred.


%Unlike full-feedback datasets used in prior routing work, which record outcomes for all models $t \in \mathcal{T}$, observational datasets contain only \emph{partial feedback}, logging the outcome of a single model $t_i$ selected by historical policies. As a result, counterfactual outcomes for unobserved models are missing, making it necessary to estimate them while correcting for treatment bias. We address this challenge using causal inference techniques, as detailed in Section~\ref{sect_DR}.


%Unlike standard full feedback datasets commonly used in existing LLM routing studies, where the accuracy and cost of every available model are recorded for each query, our observational dataset contains feedback only for the model actually selected. For a given query, we observe the performance of the chosen model but lack information about how alternative models would have performed. This partial feedback setting introduces additional challenges for evaluation and learning, as it requires reasoning about counterfactual outcomes based solely on the observed data.


%\subsection{Causal Smart Predict-and-Route Framework}\label{sec:smart}
\subsection{End-to-End Regret Minimization}\label{sec:smart}

Our goal is to route each prompt $x$ to the LLM $t$ such that the decision leads to the highest possible utility $y_x(t)$. To this end, we aim to learn an end-to-end policy $f: \mathcal{X} \to \mathcal{T}$ that minimizes the decision-making regret \citep{fernandez2022causal,pmlr-v162-zou22a}. More formally,
\begin{align*}
    f^* := \arg \min_f Regret(f)
\end{align*}
\noindent where regret is defined as
\begin{align}
    Regret(f) := \mathbb{E}_X[Y_X(t^*_X) - Y_X(f(X))] = \mathbb{E}_X[Y_X(t^*_X)] - \mathbb{E}_X[Y_X(f(X))],
\end{align}
\noindent with $t^*_X := \arg\max_{t\in\mathcal{T}} Y_X(t)$. 

We want to point out that unlike full-feedback datasets used in prior routing work, which record outcomes for all models $t \in \mathcal{T}$, observational datasets contain only \emph{partial feedback}, logging the outcome of a single model $t_i$ selected by historical policies. As a result, counterfactual outcomes for unobserved LLMs are missing, making it necessary to estimate them while correcting for treatment bias. We address this challenge of estimating $\hat{Y}_X(\cdot)$ using causal inference techniques, as detailed in Section~\ref{sect_DR}.

With an accurate approximation $\hat{Y}_X(\cdot)$,  the empirical regret can be approximated as
\begin{align} Regret(f) \approx \frac{1}{n} \sum_{i=1}^n \left( \hat{Y}_{x_i}(t^*_i) - \hat{Y}_{x_i}(f(x_i)) \right), \label{eq:approx} \end{align} 
\noindent where $t^*_i := \arg\max_{t \in \mathcal{T}} \hat{Y}_{x_i}(t)$ is the estimated optimal decision for query $x_i$. 




A key challenge with the objective in Equation~\eqref{eq:approx} is its dependence on the discrete routing decision $f(x_i)$, which makes the regret non-differentiable. To address this, we introduce two surrogate loss functions that serve as differentiable approximations.



\paragraph{Surrogate Loss 1: Classification-Based Upper Bound}


Our first approach is to derive a tractable upper bound on the regret and directly minimize it. To do so, we define the following notion of Lipschitz continuity for utility functions over the probability simplex.

\begin{definition}\label{def:lip}
Let $\hat{Y}_x : \mathcal{T} \to \mathbb{R}$ be an estimated utility function assigning a scalar utility to each model $t \in \mathcal{T}$ for a given input $x$. We extend $\hat{Y}_x$ to the probability simplex $\Delta^{|\mathcal{T}|}$ by defining
\begin{align}
\hat{Y}_x(p) := \sum_{t \in \mathcal{T}} p(t) \hat{Y}_x(t),
\end{align}
for any $p \in \Delta^{|\mathcal{T}|}$. We say that $\hat{Y}_x$ is $L$-Lipschitz over the simplex with respect to the $\ell_1$ norm if for all $p, q \in \Delta^{|\mathcal{T}|}$,
\begin{align}
|\hat{Y}_x(p) - \hat{Y}_x(q)| \leq L \cdot \|p - q\|_1.
\end{align}
Here, the constant $L$ can be taken as \( L := \max_{t \in \mathcal{T}} |\hat{Y}_x(t)| \). 
\end{definition}

This condition ensures that small changes in the model distribution lead to bounded changes in expected utility. Since $\mathcal{T}$ is a finite set and $\hat{Y}_x(\cdot)$ is typically learned via bounded smooth function approximators (e.g., neural networks), it is natural to expect bounded variation in utility values across nearby treatments. Note that when $p$ is a one-hot vector $e_{t^*}$ and $q = f(x)$ is a probabilistic policy, this setting corresponds to our model selection problem, where we seek to minimize the regret between the optimal choice and a stochastic routing decision.

\begin{proposition}
\label{prop:regret_ce_bound}
Suppose the estimated utility function $\hat{Y}_x : \mathcal{T} \to \mathbb{R}$ is $L$-Lipschitz continuous over the probability simplex with respect to the $\ell_1$ norm, as in Definition~\ref{def:lip}. Then, for a policy $f : \mathcal{X} \to \Delta^{|\mathcal{T}|}$ that outputs a distribution $f(x)$ over $\mathcal{T}$, the regret can be upper bounded by:
\begin{align}
    \text{Regret}(f) \leq L \cdot \frac{1}{n} \sum_{i=1}^n \sqrt{2 \cdot \text{CE}(t_i^*, f(x_i))},
\end{align}
where $t_i^* := \arg\max_{t \in \mathcal{T}} \hat{Y}_{x_i}(t)$ is the optimal treatment for input $x_i$, and $\text{CE}(t_i^*, f(x_i)) := -\log f(x_i)_{t_i^*}$ denotes the cross-entropy loss.
\end{proposition}



This motivates a classification-based surrogate objective: rather than modeling the full utility surface, we directly learn a policy $f : \mathcal{X} \to \mathcal{T}$ by solving a supervised learning problem, where the target label for each input $x_i$ is the estimated optimal decision $t_i^*$. Optimizing a classification loss $d(t_i^*, f(x_i))$, serves as a tractable surrogate for minimizing the regret in Equation~\eqref{eq:approx}, as it upper bounds the regret under mild assumptions. This formulation reduces policy learning to a multiclass classification task, enabling efficient training using standard techniques.


\subsubsection{Surrogate Loss 2: Softmax-Weighted Regret}



The second proxy directly minimizes the regret using a differentiable softmax approximation. Specifically, we model the policy function \( f \) as a neural network with \( |\mathcal{T}| \) outputs passed through a softmax layer with temperature parameter \( \tau > 0 \), which makes the regret surrogate in Equation~\eqref{eq:approx} differentiable. The first term of the regret, \( \mathbb{E}_X[Y_X(t^*)] \), is approximated as:
\begin{equation}
    \mathbb{E}_X[Y_X(t^*)] \approx \frac{1}{n} \sum_{i=1}^n \hat{Y}_{x_i}(t_i^*),
    \quad \text{where } t_i^* := \arg\max_{t \in \mathcal{T}} \hat{Y}_{x_i}(t).
\end{equation}
The second term, \( \mathbb{E}_X[Y_X(f(X))] \), is estimated by treating the softmax output as a distribution over treatments:
\begin{equation}
    \mathbb{E}_X[Y_X(f(X))] \approx \frac{1}{n} \sum_{i=1}^n \sum_{t=1}^{|\mathcal{T}|} 
    \hat{Y}_{x_i}(t) \cdot \mathrm{softmax}(f(x_i))_t.
\end{equation}




% The second proxy aims to directly minimize the regret using a differentiable softmax approximation. Specifically, we model the policy function $f$ as a neural network whose $|\mathcal{T}|$ outputs are passed through a softmax layer with temperature parameter $\tau > 0$, rendering the regret approximation in Equation~\eqref{eq:approx} differentiable. The first term of the regret, $\mathbb{E}_X[Y_X(t^*)]$, can be approximated as:
% \begin{align}
%     \mathbb{E}_X[Y_X(t^*)] \approx \frac{1}{n} \sum_{i=1}^n \hat{Y}_{x_i}(t_i^*),
% \end{align}
% \noindent where $t_i^* := \arg\max_{t \in \mathcal{T}} \hat{Y}_{x_i}(t)$ denotes the estimated optimal treatment for input $x_i$. To approximate the second term $\mathbb{E}_X[Y_X(f(X))]$, we treat the softmax output as a probability distribution over treatments:
% \begin{align}
%     \mathbb{E}_X[Y_X(f(X))] \approx \frac{1}{n} \sum_{i=1}^n \sum_{t=1}^{|\mathcal{T}|} \hat{Y}_{x_i}(t) \cdot \mathrm{softmax}(f(x_i))_t.
% \end{align}

Combining the two, we minimize the following differentiable surrogate objective:
\begin{align}
    \min_f \ \frac{1}{n} \sum_{i=1}^n \Bigg( \hat{Y}_{x_i}(t_i^*) - \sum_{t=1}^{|\mathcal{T}|} \hat{Y}_{x_i}(t) \cdot \mathrm{softmax}(f(x_i))_t \Bigg).
\end{align}

After training, the learned policy prescribes for each $x \in \mathcal{X}$ the treatment $ \hat{t} = \arg\max_{t \in \mathcal{T}} f(x)_t.$ 
We now show that this objective recovers pointwise optimal treatment assignment, thus providing a consistent and differentiable approximation to the original regret minimization objective. %minimizes the regret approximation at convergence.

\begin{proposition}
Let $f: \mathcal{X} \to \mathbb{R}^{|\mathcal{T}|}$ be a neural network whose output is passed through a softmax layer with fixed temperature $\tau > 0$, and define $t^*_i := \arg\max_{t \in \mathcal{T}} \hat{Y}_{x_i}(t)$. Then, optimizing the softmax-weighted surrogate regret objective via gradient descent 
\begin{align}
    \min_f \frac{1}{n} \sum_{i=1}^n \bigg( \hat{Y}_{x_i}(t^*_i) - \sum_{t=1}^{|\mathcal{T}|} \hat{Y}_{x_i}(t) \cdot \mathrm{softmax}(f(x_i))_t \bigg)
\end{align}
leads the model $f$ to concentrate all probability mass on the optimal treatment $t^*_i$. That is, at convergence,\begin{equation}
    \mathrm{softmax}(f(x_i))_t \to 
    \begin{cases}
        1 & \text{if } t = t^*_i, \\
        0 & \text{otherwise}.
    \end{cases}
\end{equation}
\end{proposition}

%The result shows that optimizing the softmax-weighted surrogate leads to a policy that recovers pointwise optimal treatment assignment, thus providing a consistent and differentiable approximation to the original regret minimization objective.

% \begin{proof}
% Let $\hat{Y}_{x_i} \in \mathbb{R}^{|\mathcal{T}|}$ denote the vector of estimated potential outcomes for input $x_i$, and let $f(x_i) \in \mathbb{R}^{|\mathcal{T}|}$ be the output of the neural network before the softmax layer. The objective for a single instance $x_i$ can be written as minimizing the regret surrogate:
% \begin{align}
%     \hat{Y}_{x_i}(t^*_i) - \sum_{t=1}^{|\mathcal{T}|} \hat{Y}_{x_i}(t) \cdot \mathrm{softmax}(f(x_i))_t.
% \end{align}
% This is equivalent to maximizing the inner product:
% \begin{align}
%     \langle \hat{Y}_{x_i}, \mathrm{softmax}(f(x_i)) \rangle.
% \end{align}

% Let us denote $p := \mathrm{softmax}(f(x_i)) \in \Delta^{|\mathcal{T}| - 1}$, the probability simplex. We now show that the inner product $\langle \hat{Y}_{x_i}, p \rangle$ increases at each gradient step. Since $p = \mathrm{softmax}(f(x_i))$, we can compute the gradient of the objective with respect to $f(x_i)$ as:
% \begin{align}
%     \nabla_{f(x_i)} \langle \hat{Y}_{x_i}, \mathrm{softmax}(f(x_i)) \rangle
%     = J_{\mathrm{softmax}}(f(x_i))^\top \hat{Y}_{x_i},
% \end{align}
% where $J_{\mathrm{softmax}}(f(x_i))$ is the Jacobian of the softmax function, given by:
% \begin{align}
%     J_{\mathrm{softmax}}(f(x_i))_{t,s} = \frac{\partial \, \mathrm{softmax}(f(x_i))_t}{\partial f(x_i)_s}
%     = \mathrm{softmax}(f(x_i))_t \left(\delta_{t,s} - \mathrm{softmax}(f(x_i))_s \right).
% \end{align}

% This gradient direction corresponds to increasing the logit value of actions with higher $\hat{Y}_{x_i}(t)$ and decreasing those with lower values, pushing the softmax distribution toward the mode of $\hat{Y}_{x_i}$. In other words, the gradient ascent step increases the inner product at each iteration $k$:
% \begin{align}
%     \langle \hat{Y}_{x_i}, \mathrm{softmax}(f(x_i)) \rangle^{(k+1)} 
%     > \langle \hat{Y}_{x_i}, \mathrm{softmax}(f(x_i)) \rangle^{(k)}.
% \end{align}

% Since $\hat{Y}_{x_i}$ is fixed and the softmax is smooth and bounded, this sequence is monotonically increasing and converges to the maximum possible value:
% \begin{align}
%     \langle \hat{Y}_{x_i}, \mathrm{softmax}(f(x_i)) \rangle \to \max_t \hat{Y}_{x_i}(t) = \hat{Y}_{x_i}(t^*_i),
% \end{align}
% which implies:
% \begin{align}
%     \mathrm{softmax}(f(x_i))_t \to 
%     \begin{cases}
%         1 & \text{if } t = t^*_i, \\
%         0 & \text{otherwise}.
%     \end{cases}
% \end{align}

% Thus, the regret surrogate converges to zero:
% \begin{align}
%     \hat{Y}_{x_i}(t^*_i) - \langle \hat{Y}_{x_i}, \mathrm{softmax}(f(x_i)) \rangle \to 0,
% \end{align}
% and the learned policy selects the treatment maximizing the estimated outcome.
% \end{proof}

\subsection{Estimating Counterfactual Utility via Causal Inference}\label{sect_DR}

In the previous sections, we assumed that we have access to $\hat{Y}_X(\cdot)$. Given the observational nature of the data, the potential utility function $Y_x(\cdot)$ is not directly observable. We follow the potential outcomes framework \cite{rosenbaum1983central, rubin1984bayesianly} and assume the existence of a potential utility function $Y_x(t)$. %Our goal is to learn a single end-to-end prescriptive model that, given a prompt $x \in \mathcal{X}$, routes it to the LLM $t \in \mathcal{T}$ that maximizes the utility $y_x(t) \in \mathcal{Y}$. 
We adopt the following assumptions, which are standard in the causal inference literature.


\begin{assumption}[Stable Unit Treatment Value]
The potential outcome of one sample is independent of the treatment assignments on the other samples.
\end{assumption}

\begin{assumption}[Ignorability]
The assigned treatments and potential outcomes are independent conditional on observed covariates, i.e. $t\perp \{Y_x(t')|t'\in \mathcal{T}\}|x$ \text{\citep{hirano2004propensity}}.
\end{assumption}


\begin{assumption}[Support]
For $x\in \mathcal{X}$ such that $p(x)>0$, we have that $p(t|x)>0$ for each $t\in\mathcal{T}$.
\end{assumption}

In the causal inference literature, counterfactual outcomes can be estimated using various methods. Such examples include the  ``meta-learner''~\citep{kunzel2019metalearners}, or the Inverse Propensity Weighting (IPW) estimator~\citep{horvitz1952generalization}. In this work, we utilize the doubly robust estimator introduced by \cite{dudik2011doubly}, defined as: 
\begin{align} \hat{Y}_x(t) := \frac{(y - \hat{r}_{t}(x)) \mathbbm{1}[\pi(x) = t]}{\hat{p}(t|x)} + \hat{r}_{t}(x), \quad \forall t \in \mathcal{T}, \end{align} 
\noindent where $\hat{r}_{t} : \mathcal{X} \to \mathcal{Y}$ denotes the direct outcome regression model for treatment $t$, $\hat{p}(t|x)$ is the estimated propensity score, and $\pi : \mathcal{X} \to \mathcal{T}$ is the logging policy observed in the dataset $\mathcal{D}$.






%The doubly robust estimator is particularly well-suited to observational settings, as it leverages two complementary sources of information—the outcome model $\hat{r}_t(x)$ and the propensity model $\hat{p}(t|x)$, to mitigate bias in the estimation of counterfactual outcomes. It is termed doubly robust because it provides consistent estimates as long as either the outcome model or the propensity model is correctly specified, but not necessarily both. \citet{dudik2011doubly} demonstrated that this estimator achieves a favorable bias–variance trade-off: the propensity score corrects for selection bias, while the outcome model reduces variance by incorporating structural information about the utility function.



The doubly robust estimator combines an outcome model $\hat{r}_t(x)$ and a propensity model $\hat{p}(t|x)$, yielding consistent estimates if either is correctly specified~\citep{dudik2011doubly}. It offers a favorable bias-variance trade-off: the propensity model corrects for selection bias, while the outcome model reduces variance by leveraging structure in the data.\\


\begin{remark}
     While we use the doubly robust (DR) estimator in our experiments due to its favorable bias–variance tradeoff and strong practical performance, our framework is estimator-agnostic: any valid counterfactual estimator can be used to compute utility estimates. This includes more advanced approaches that relax or mitigate these assumptions
\end{remark}
    


In the experimental section, we show that %accounting for treatment bias in observational data significantly improves performance. 
ignoring the treatment bias leads to inaccurate counterfactual estimates and causes substantial degradation in routing quality, highlighting the limitations of standard supervised learning approaches that assume full feedback.


\section{Routing under Heterogeneous Cost Preferences}\label{sec:multi_task}





In the previous section, we introduced a causal end-to-end framework for learning optimal routing policies from observational data, where the objective is to maximize a utility function of the form $y = a - \lambda c$, with a \emph{fixed} \( \lambda \geq 0 \) representing the trade-off between accuracy and cost.

In practice, however, user preferences vary, i.e., different queries may be associated with different values of $\lambda$. From a system design perspective, training and maintaining a separate router for each possible $\lambda$ is impractical. In this section, we propose a unified approach that supports routing under heterogeneous cost sensitivities. We first present a joint model architecture that conditions on both the query and the cost parameter, and then provide a theoretical analysis to justify the proposed design.



\subsection{Interval-Conditioned Joint Router}\label{sec:model_Setup}

%To address this limitation, we propose a \emph{multi-task, budget-aware router} that generalizes across a set of possible budget constraints. 
We design a neural network $f : \mathcal{X} \times \mathbb{R}_{\geq 0} \to \mathbb{R}^{|\mathcal{T}|}$ that jointly takes a query $x \in \mathcal{X}$ and a cost sensitivity parameter $\lambda \in \mathbb{R}_{\geq 0}$ as input, outputs a score vector over available LLMs and the routing decision is then made via:\begin{align}
    \hat{t} = \arg\max_{t \in \mathcal{T}} f(x, \lambda)_t.
\end{align}
We assume access to a finite, representative set of cost preferences $\Lambda := \{\lambda_1, \dots, \lambda_m\} \subset \mathbb{R}_{\geq 0}$. For ease of notation we assume that $\lambda_1 < \lambda_2 < \dots < \lambda_m$. In practice, these values may correspond to discrete service quality tiers (e.g., basic, standard, premium) that reflect users’ varying willingness to trade off cost for performance.  
For each $\lambda \in \Lambda$, we partition the training data by cost preference and estimate $\lambda$-specific utility $\hat{Y}_{x_i}^{\lambda}(t)$, which forms the basis of our joint interval-conditioned architecture.




\paragraph{Training Procedure.} %We adopt a hybrid strategy:
\begin{enumerate}[leftmargin=1.5em]
    \item %\textbf{(Train Per-$\lambda$ Models) }
    For each $\lambda \in \Lambda$, we first train an individual router $f_\lambda : \mathcal{X} \to \mathbb{R}^{|\mathcal{T}|}$ using the methods introduced in Section~\ref{sec:smart}.
    \item %\textbf{(Joint Routing Model Initialization)} 
    For each interval $[\lambda_j, \lambda_{j+1}]\, j \in [1,\dots,m-1]$,  we initialize a joint network $f(x, \lambda): \mathcal{X} \times \mathbb{R}_{\geq0} \to \mathbb{R}^{|\mathcal{T}|}$ that uses as input both $x$, $\lambda \in (\lambda_j, \lambda_{j+1})$.
    \item The shared model is fine-tuned to minimize regret over the interval:
    \begin{align}
        \min_f \ \frac{1}{2n} \sum_{\lambda \in \{\lambda_j, \lambda_{j+1}\}} \sum_{i=1}^n \text{Regret}\left(f(x_i, \lambda)\right),
    \end{align}
    where regret is computed using the doubly robust estimator as described earlier, under the corresponding $\lambda$-specific utility $\hat{Y}_{x_i}^{\lambda}(t)$.
\end{enumerate}

\paragraph{Deployment Strategy.} At inference time, given a user-specified cost sensitivity parameter $\lambda \in \mathbb{R}_{\geq 0}$:
\begin{itemize}[leftmargin=1.5em]
    \item If $\lambda \in \Lambda$, we use the individual model $f_{\lambda}(x)$.
    \item If $\lambda \notin \Lambda$, we identify the closest neighbors $\underline{\lambda}, \overline{\lambda} \in \Lambda$ such that $\underline{\lambda} < \lambda < \overline{\lambda}$, and we use the corresponding joint network $f(x, \lambda)$ trained to generalize across the preference in $(\underline{\lambda}, \overline{\lambda} )$.
\end{itemize}

%This formulation provides a unified and flexible solution that avoids retraining for each new budget constraint.

\subsection{Model Architecture}

A key component of the heterogeneous cost preference routing setup introduced in Section~\ref{sec:model_Setup} is the interval-conditioned joint model $f(x, \lambda)$, that for each interval $[\lambda_j, \lambda_{j+1}]$, interpolates between the two corresponding individual models $f_{\lambda_j}$ and $f_{\lambda_{j+1}}$.
Specifically, the architecture is designed to exploit the \emph{affine structure} of the utility function with respect to $\lambda$, namely $y = a - \lambda c$. This motivates a lightweight parameterization that uses only the two endpoints of the interval $[\lambda_j, \lambda_{j+1}]$ rather than all $m$ pre-trained models. Concretely, the joint model that we propose is defined as:
\begin{align}
    f(x, \lambda) = \texttt{Linear}\left( [f_{\lambda_j}(x),\ f_{\lambda_{j+1}}(x)]+\ g(\lambda) \right),
\end{align}
where $[\cdot,\cdot]$ denotes concatenation, and $g(\lambda) := \texttt{Activation}(\texttt{Linear}(\lambda))$ is a learnable representation of the cost sensitivity parameter.

This architecture enables smooth interpolation between $f_{\lambda_j}$ and $f_{\lambda_{j+1}}$ within $[\lambda_j, \lambda_{j+1}]$, allowing the router to adapt to intermediate values of $\lambda$ without requiring an individual model for each one. By conditioning only on the two bounding models, this design achieves computational efficiency and strong generalization across cost preferences. The proposed architecture is illustrated in Figure~\ref{fig:budget-aware-architecture}.


\begin{figure}[h]
\centering
\includegraphics[width=0.82\linewidth]{images/routing_arc_3.png}
\caption{Overview of the proposed interval-conditioned joint router framework. \textbf{Left:} Decision logic for handling a given cost sensitivity parameter $\lambda$. \textbf{Right:} Joint router architecture.}
\label{fig:budget-aware-architecture}
\end{figure}

\subsection{Theoretical Guarantees}

We now present theoretical guarantees that justify the structure and training strategy of joint cost-preference router. These guarantees leverage the affine nature of the utility function, %the normalized accuracy objective $y = a - \lambda \cdot c$, 
which is linear in the cost parameter $\lambda$ for fixed accuracy $a$ and cost $c$, enabling exact interpolation for specific cost sensitivity across fixed  intervals.

\begin{proposition}[Piecewise Constant Optimal Policy]
\label{prop:piecewise-policy}
Fix a query $x \in \mathcal{X}$ and assume the estimated utility is affine in $\lambda$, i.e., $\hat{Y}_x^\lambda(t) = a_x(t) - \lambda \cdot c_x(t)$ for all $t \in \mathcal{T}$. Then the optimal treatment
\[
t^*(\lambda) := \arg\max_{t \in \mathcal{T}} \hat{Y}_x^\lambda(t)
\]
is piecewise constant in $\lambda$. That is, the cost parameter $\mathbb{R}_{\geq 0}$ can be partitioned into intervals over which the optimal treatment remains fixed.
\end{proposition}

% \begin{proof}
% For fixed $x$, each $\hat{Y}_x^\lambda(t)$ is an affine function of $\lambda$. The pointwise maximum of a finite collection of affine functions is piecewise affine, and the argmax corresponds to the highest line at each $\lambda$. Since each pair of lines can intersect at most once, the number of intervals over which a single treatment is optimal is bounded by $|\mathcal{T}| - 1$. Therefore, $t^*(\lambda)$ changes only at these intersection points and remains constant within each interval.
% \end{proof}

\begin{proposition}[Affine Closure of Utility Function]
\label{prop:affine-closure}
Let $\lambda_j < \lambda_{j+1}$ be two adjacent cost values and let $\lambda \in [\lambda_j, \lambda_{j+1}]$. Suppose the utility function is affine in $\lambda$, i.e., $\hat{Y}_x^\lambda(t) = a_x(t) - \lambda \cdot c_x(t).$ Then for all $t \in \mathcal{T}$, the utility at $\lambda$ is a convex combination of utilities at the endpoints:
\[
\hat{Y}_x^\lambda(t) = \alpha \cdot \hat{Y}_x^{\lambda_j}(t) + (1 - \alpha) \cdot \hat{Y}_x^{\lambda_{j+1}}(t),
\quad \text{where } \alpha := \frac{\lambda_{j+1} - \lambda}{\lambda_{j+1} - \lambda_j}.
\]
\end{proposition}

% \begin{proof}
% We expand each term:
% \begin{align*}
% \hat{Y}_x^{\lambda_j}(t) &= a_x(t) - \lambda_j \cdot c_x(t), \\
% \hat{Y}_x^{\lambda_{j+1}}(t) &= a_x(t) - \lambda_{j+1} \cdot c_x(t).
% \end{align*}
% Then:
% \begin{align*}
% \alpha \cdot \hat{Y}_x^{\lambda_j}(t) + (1 - \alpha) \cdot \hat{Y}_x^{\lambda_{j+1}}(t)
% &= \alpha \cdot (a_x(t) - \lambda_j c_x(t)) + (1 - \alpha) \cdot (a_x(t) - \lambda_{j+1} c_x(t)) \\
% &= a_x(t) - [\alpha \lambda_j + (1 - \alpha) \lambda_{j+1}] \cdot c_x(t) \\
% &= a_x(t) - \lambda \cdot c_x(t) = \hat{Y}_x^\lambda(t),
% \end{align*}
% since:
% \[
% \alpha \lambda_j + (1 - \alpha) \lambda_{j+1} = \lambda.
% \]
% \end{proof}

\begin{corollary}[Sufficiency of Two Models per Interval]
\label{cor:two-models}
Under the affine assumption, the utility $\hat{Y}_x^\lambda(t)$ for any $\lambda \in [\lambda_j, \lambda_{j+1}]$ can be exactly reconstructed using only the endpoints $\hat{Y}_x^{\lambda_j}(t)$ and $\hat{Y}_x^{\lambda_{j+1}}(t)$. Thus, it is sufficient to use only the two corresponding models $f_{\lambda_j}$ and $f_{\lambda_{j+1}}$ for interpolation within the interval.
\end{corollary}
\begin{proposition}[Expressivity of Additive Two-Model joint Architecture]
\label{prop:architecture-expressivity}
Let $\lambda \in [\lambda_j, \lambda_{j+1}]$, and suppose that for each $t \in \mathcal{T}$ the utility satisfies $\hat{Y}_x^\lambda(t) = a_x(t) - \lambda \cdot c_x(t)$. Then the optimal treatment $t^*(\lambda) := \arg\max_{t} \hat{Y}_x^\lambda(t)$ can be exactly represented by a softmax policy over a function of the form:
\[
f(x, \lambda) = \texttt{Linear} \left( [f_{\lambda_j}(x), f_{\lambda_{j+1}}(x)] + g(\lambda) \right),
\]
where $g(\lambda)$ is any differentiable embedding of $\lambda$, and $f_{\lambda_j}, f_{\lambda_{j+1}}$ are accurate predictors trained at endpoints $\lambda_j$ and $\lambda_{j+1}$.
\end{proposition}


% \begin{proof}
% From Proposition~\ref{prop:affine-closure}, the utility $\hat{Y}_x^\lambda(t)$ is a convex combination of $\hat{Y}_x^{\lambda_j}(t)$ and $\hat{Y}_x^{\lambda_{j+1}}(t)$. If the network $f(x, \lambda)$ linearly combines the outputs of $f_{\lambda_j}(x)$ and $f_{\lambda_{j+1}}(x)$, then its scores can match $\hat{Y}_x^\lambda(t)$ up to a scalar transformation. Applying softmax preserves the argmax.

% Including $g(\lambda)$ allows the architecture to learn any additional monotonic reweighting of the interpolation, ensuring the output scores can be shaped to approximate the true utility surface exactly. Thus, the architecture can represent the optimal policy within each interval.
% \end{proof}


\paragraph{Implications for Architecture and Training.}
The theoretical results above provide strong justification for both the proposed model architecture and the associated training procedure. Proposition~\ref{prop:piecewise-policy} shows that the optimal treatment changes only across a small number of cost sensitivities, supporting our interval-conditioned strategy for routing. Proposition~\ref{prop:affine-closure} guarantees that the utility for any intermediate cost preference can be exactly recovered through convex interpolation of endpoint models. Finally, Proposition~\ref{prop:architecture-expressivity} establishes that our joint model  is expressive enough to capture the optimal policy within each interval. 
%In practice, the set of representative cost parameters $\Lambda$ can be seen as a menu of cost-performance trade-offs offered to users (e.g., a pricing or quality-of-service tier). Alternatively, users may directly specify their own cost sensitivity. 
Our method provides a principled approach for learning a joint routing model that accommodates heterogeneous cost preferences from observational data.


\section{Experiments}\label{sec:exp}

%We evaluate the performance of our proposed methods on two real-world datasets, comparing them against a range of competitive baselines. 
% All experiments were implemented in Python 3.8.12~\cite{python}, using PyTorch 2.4.1+cu121~\cite{NEURIPS2019_9015} and Scikit-learn~\cite{sklearn_api}. Experiments were conducted on an internal compute cluster equipped with an Intel(R) Xeon(R) Platinum 8260 CPU @ 2.40GHz, 512~GB of RAM, and two NVIDIA V100 GPUs with 16~GB memory each.


\subsection{Datasets}

We evaluate our methods on two publicly available benchmarks for LLM routing: \textbf{RouterBench}~\citep{hu2024routerbench} and \textbf{SPROUT}~\citep{somerstep2025carrot}.

\textbf{RouterBench} is a standardized benchmark comprising 35{,}712 prompt–response pairs from 11 language models. The prompts are drawn from eight evaluation suites spanning reasoning, factual recall, dialogue, mathematics, and code generation. Each prompt is annotated with model accuracy and execution cost, enabling supervised training and evaluation of routing policies.

\textbf{SPROUT} is a larger and more diverse benchmark focused on cost-aware routing, consisting of 44{,}241 prompts and responses from 13 state-of-the-art LLMs. The prompts cover six challenging tasks: GPQA~\citep{rein2024gpqa}, MuSR~\citep{sprague2023musr}, MMLU-Pro~\citep{wang2024mmlu}, MATH~\citep{hendrycks2021measuring}, OpenHermes~\citep{teknium2023openhermes}, and RAGBench~\citep{friel2407ragbench}. SPROUT includes a predefined split: 80\% for training, with the remaining 20\% evenly divided between validation and test sets.

% We conduct experiments on two publicly available benchmark datasets for LLM routing: \textbf{RouterBench}~\citep{hu2024routerbench} and \textbf{SPROUT}~\citep{somerstep2025carrot}.

% \textbf{RouterBench} is a standardized benchmark for LLM routing, comprising 35,712 prompt-response pairs collected from 11 LLMs. The prompts span eight different evaluation benchmarks covering reasoning, factual knowledge, dialogue, mathematics, and code generation. Each prompt is annotated with model accuracy and execution cost, enabling response-based decision-making. 

% \textbf{SPROUT} is a more recent and larger benchmark for cost-aware routing, consisting of 44,241 prompts and responses from 13 state-of-the-art language models. The prompts are drawn from six diverse benchmarks, including GPQA \citep{rein2024gpqa}, MuSR \citep{sprague2023musr}, MMLU-Pro \citep{wang2024mmlu}, MATH \citep{hendrycks2021measuring}, OpenHermes \citep{teknium2023openhermes}, and RAGBench \citep{friel2407ragbench}. SPROUT includes a predefined train/validation/test split, using 80\% of the data for training and splitting the remaining 20\% equally between validation and test sets.




\subsection{Embeddings and Model Architecture}

% To encode input queries into vector representations \(x\), we employ a two-stage embedding process. First, we enrich each prompt with contextual metadata by prepending a natural language prefix that identifies the source dataset. Specifically, for a prompt \texttt{p} originating from dataset \texttt{D} (e.g., \texttt{openhermes/teknium}), we construct the following context-augmented input: \textit{``The following prompt comes from the dataset D. The prompt is: \texttt{p}''}. This step provides the embedding model with useful dataset-level context, which is particularly beneficial in multi-domain routing scenarios. The template is flexible and can be extended to include additional metadata if desired.

% To encode input queries into vector representations \(x\) we generate embeddings using two lightweight, publicly available language models: \texttt{BERT-base-uncased} (embed size 768) and \texttt{Llama-3.2-1B} (embed size 2048). These models were selected due to their small size and efficient inference, making them suitable for real-time routing applications. For each prompt, we perform a single feedforward pass through the model and apply mean pooling over the final hidden states to produce a fixed-dimensional sentence embedding. 

% The resulting embeddings are then fed into a two-layer fully connected neural network with GELU activations and 200 hidden units per layer. The model is trained for a maximum of 10{,}000 epochs using the Adam optimizer with a learning rate of \(1\mathrm{e}{-4}\), and early stopping with a patience of 100 epochs. The model uses a softmax output layer with a temperature of \(\tau = 100\), which controls the sharpness of the output probabilities. This model architecture is used consistently across all benchmarked methods to ensure a fair comparison. For the doubly robust estimator, we use the same architecture for learning the direct outcome models \( \hat{r}_t(x) \), while employing an XGBoost model to estimate the propensity scores \( \hat{p}(t \mid x) \). To mitigate instability from extreme inverse propensity weights, we apply clipping at the 5th and 95th percentiles. The hyperparameters used are reported in Appendix~\ref{app:hyperparams}. The only difference is that for the interval architecture, we increase the softmax temperature to \(\tau = 1000\) to better support smooth interpolation across different \(\lambda\). 

To encode input queries into vector representations \(x\), we generate embeddings using two compact, publicly available language models: \texttt{BERT-base-uncased} (768 dimensions) and \texttt{Llama-3.2-1B} (2048 dimensions). Each input is passed once through the model, and the final hidden states are mean-pooled to obtain a fixed-length embeddings. These models were selected for their efficiency and suitability for real-time routing.

The embeddings are processed by a two-layer fully connected network with GELU activations and 200 hidden units per layer. The model is trained with the Adam optimizer (learning rate \(1 \times 10^{-4}\)) for up to 10{,}000 epochs, using early stopping with a patience of 100. A softmax output with temperature \(\tau = 100\) is used to control the sharpness of the output probabilities. This architecture is used consistently across all benchmarked methods for fair comparison. For the doubly robust estimator, the same network models the direct outcomes \(\hat{r}_t(x)\), while the propensity scores \(\hat{p}(t \mid x)\) are estimated using XGBoost. To reduce variance from extreme inverse propensity weights, we apply clipping at the 5th and 95th percentiles. The only architectural modification is for the interval-based model, where the softmax temperature is increased to \(\tau = 1000\) to enable smoother interpolation across \(\lambda\). Hyperparameters are summarized in Appendix~\ref{app:hyperparams}.






% This embedding and model architecture yields semantically meaningful, task-aware representations that are compatible with a wide range of routing algorithms, while remaining computationally efficient and scalable across datasets.


\subsection{Methods}

% We evaluate our proposed routing strategies against a broad range of benchmarks from the causal machine learning and LLM routing literature. Since both SPROUT and RouterBench provide full-feedback datasets (i.e., responses from all candidate models), we simulate observational data by sampling a single treatment per prompt. Specifically, for each prompt, we select a model \( t \in \mathcal{T} \) with probability proportional to its accuracy via $\mathbb{P}[t = \tau] = \frac{e^{a_{\tau}}}{\sum_{\tau' \in \mathcal{T}} e^{a_{\tau'}}},$ where \( a_{\tau} \) denotes the accuracy of model \( \tau \) on the given prompt. 


% \paragraph{Full-Feedback (Upper Bound).}  
% We first train a routing model \( f: \mathcal{X} \to \mathbb{R}^{|\mathcal{T}|} \) using the full-feedback dataset under a standard multi-class classification loss. Since this model observes all potential outcomes during training, it serves as an optimistic upper bound on performance.

% \paragraph{Direct Prediction Decoupled (Baseline 1).}  
% We include a baseline that takes as input the assigned model \( t \) and prompt embedding \( x \), and trains two regressors \( f: \mathcal{X} \times \mathcal{T} \to \mathbb{R} \) to predict the accuracy and the cost separately. This model is trained using standard supervised learning, without correcting for treatment assignment bias.

% \paragraph{Direct Prediction (Baseline 2).}  
% We include a baseline that takes as input the assigned model \( t \) and prompt embedding \( x \), and trains a regressor \( f: \mathcal{X} \times \mathcal{T} \to \mathbb{R} \) to predict the observed utility $y=a-\lambda \cdot c$. This model is trained using standard supervised learning, without correcting for treatment assignment bias.


% \paragraph{Reweighting via Regress-and-Compare (R\&C).}  
% To address selection bias in the observational setting, we apply the Regress-and-Compare (R\&C) technique. For each LLM \( t \), we fit a separate outcome regression model \( \hat{r}_t(x) \), then use the predicted outcomes across all \( t \in \mathcal{T} \) to identify the best model (highest estimated response) for each query. This method acts as a plug-in policy based on outcome estimation.

% \paragraph{CARROT.}  
% We evaluate the CARROT router~\cite{somerstep2025carrot}, which includes both a parametric model trained on sentence embeddings and a non-parametric kNN variant. CARROT has been shown to outperform prior baselines such as RouteLLM~\cite{ong2406routellm} and RoRF~\cite{jain2023rorf}. To adapt CARROT to the observational setting, we apply the R\&C framework to its output predictions.

% \paragraph{Regression}  
% We use a sequential approach where we model $f$ as a function $f: \mathcal{X} \to \mathbb{R}^{|\mathcal{T}|}$ and solve the following squared loss minimization problem: $\min_f \ \sum_{i=1}^n \sum_{t=1}^{|\mathcal{T}|}  (\hat{Y}_{x_i}(t) - f(x_i)_t )^2$. After training $f$, we select the treatment for each input via $\hat{t} = \arg\max_{t \in \mathcal{T}} f(x)_t$.

% \paragraph{Ours.}  
% Finally, we evaluate the routing methods introduced in this work, including classification-based and regret-minimizing variants, as described in Section~\ref{sec:smart}. For the budget-aware setting, we evaluate the interpolation approach presented in Section~\ref{sec:multi_task}, which supports continuous generalization across user-defined trade-off parameters \( \lambda \).


We evaluate our proposed routing strategies against a range of baselines from the causal machine learning and LLM routing literature. Since both SPROUT and RouterBench provide full-feedback datasets (i.e., responses from all models), we simulate observational data by sampling a single model per prompt. Specifically, for each prompt, we sample a model \( t \in \mathcal{T} \) with probability proportional to its accuracy $\mathbb{P}[t = \tau] = \frac{e^{a_{\tau}}}{\sum_{\tau' \in \mathcal{T}} e^{a_{\tau'}}},$ where \( a_{\tau} \) is the accuracy of model \( \tau \) on that prompt.

As an optimistic oracle, we include a \textbf{Full-Feedback} model that learns a model \( f: \mathcal{X} \to \mathbb{R}^{|\mathcal{T}|} \) using the complete outcome vector for each query and optimizes a standard multi-class classification loss. We benchmark against a common decoupled routing strategy denoted \textbf{Baseline}, which independently estimates model accuracy \( a_x(t) \) and cost \( c_x(t) \), without accounting for selection bias. This reflects the approach taken in prior predictive routing methods such as {CARROT}~\citep{somerstep2025carrot}, which has demonstrated superior performance over alternatives like RouteLLM~\citep{ong2406routellm} and RoRF~\citep{jain2023rorf}. To adjust for treatment assignment bias, we consider a \textbf{Regress-and-Compare (R\&C)} method, which fits outcome models \( \hat{Y}_x(t) \) for each treatment \( t \), and selects the action \( \hat{t} = \arg\max_{t} \hat{Y}_x(t) \). Building on this, we implement a \textbf{Causal-CARROT} variant by adapting both the parametric and kNN instantiations of CARROT to the R\&C framework. We additionally include \textbf{CF-Regression}, which models \( f: \mathcal{X} \to \mathbb{R}^{|\mathcal{T}|} \) and is trained to minimize MSE against the counterfactual utility function from the doubly robust estimator: $\min_f \ \sum_{i=1}^n \sum_{t=1}^{|\mathcal{T}|} ( \hat{Y}_{x_i}(t) - f(x_i)_t )^2.$ Decisions are made by selecting the treatment with the highest predicted value, i.e., \( \hat{t} = \arg\max_t f(x)_t \). 

%Finally, we evaluate our proposed regret-minimization methods. \textbf{RM-Classification} serves as a classification-based upper bound that reframes the regret minimization problem as multi-class prediction over optimal treatments. In contrast, \textbf{RM-Softmax} directly optimizes a differentiable surrogate of the regret via a softmax-weighted formulation over treatment utilities. In the heterogeneous cost preference setting, we additionally evaluate \textbf{RM-Interval} router (Section~\ref{sec:multi_task}), which generalizes across a continuum of cost-sensitivity parameters \( \lambda \) by interpolating between models trained at discrete preference endpoints.

% Finally, we evaluate our regret-minimization methods. \textbf{RM-Classification} reframes regret minimization as a multi-class prediction task over optimal treatments, serving as a classification-based upper bound. \textbf{RM-Softmax} directly optimizes a differentiable surrogate using a softmax-weighted utility formulation. In the heterogeneous preference setting, we also assess \textbf{RM-Interval} (Section~\ref{sec:multi_task}), which generalizes across a continuum of cost sensitivities by interpolating between models trained at discrete \( \lambda \) values.

Finally, we evaluate our regret-minimization methods. \textbf{RM-Classification} formulates the task as multi-class prediction over optimal treatments, serving as a classification-based upper bound. \textbf{RM-Softmax} directly minimizes a softmax-weighted regret surrogate. In the heterogeneous preference setting, we also assess \textbf{RM-Interval} (Section~\ref{sec:multi_task}), which generalizes across a continuum of cost sensitivities by interpolating between models trained at discrete \( \lambda \) values.





\begin{table}[h]
\centering
\caption{Comparison of routing methods by causal reasoning and end-to-end training.}
\label{tab:causal_end2end}
\begin{tabular}{lcc}
\toprule
\textbf{Method} & \textbf{Causal} & \textbf{End-to-End} \\
\midrule
Full-Feedback     & \cmark & \cmark \\ \hdashline
Baseline      & \xmark & \xmark \\
% Baseline 2      & \xmark & \xmark \\
R\&C  & \cmark & \xmark \\
 
% CARROT (Embed)        & \cmark & \xmark \\
Causal-CARROT (kNN \& Embed)             & \cmark & \xmark \\
CF-Regression         & \cmark & \xmark \\ \hdashline
RM-Classification (Ours)    & \cmark & \cmark \\
RM-Softmax (Ours)            & \cmark & \cmark \\
RM-Interval  (Ours)     & \cmark & \cmark \\
\bottomrule
\end{tabular}
\end{table}
\subsection{Evaluations}

%We present experimental results for both the fixed-\(\lambda\) and the multi-task budget-aware routing setup. 

%\paragraph{Evaluation Protocol.}
% We evaluate our methods in two settings. In the $\lambda$-specific setting, we treat each cost sensitivity parameter \(\lambda \in \{0, 100, 200, \dots, 1000\}\) as a separate routing problem, training and evaluating each router for each $\lambda$. In the heterogeneous cost preference setting, we train \textbf{RM-Interval} on a subset of values \(\{0, 200, 400, \dots, 1000\}\) and evaluate its generalization to the held-out values \(\{100, 300, 500, 700, 900\}\). 

We evaluate our methods in two settings. In the \(\lambda\)-specific setting, each cost sensitivity value \(\lambda \in \{0, 100, \dots, 1000\}\) defines a separate routing task, with models trained and evaluated independently. In the heterogeneous preference setting, \textbf{RM-Interval} is trained on a subset \(\{0, 200, \dots, 1000\}\) and tested on held-out values \(\{100, 300, 500, 700, 900\}\) to assess generalization across cost sensitivities.



%\paragraph{Evaluation Metric.}  
%We evaluate each routing policy by reporting the achieved utility $y$ per router, which captures the trade-off between model performance and computational cost.

We report the average utility across 10 independent trials for each routing method for SPROUT dataset and BERT embeddings in Table~\ref{tab:sprout_BERT_main}, where each trial involves randomly sampling observational data and retraining all models. In Figure~\ref{fig:BERT_sprout}, we also visualize the corresponding accuracy–cost curve associated with each method. Additional plots for the rest of the datasets and embeddings, along with detailed router performance %and the AUC achieved by each router, 
are provided in Appendix~\ref{app:single_lambda_results}.


% \begin{table}[h]
% \centering
% \caption{Utility on \textbf{SPROUT} with \texttt{BERT} embeddings. \textit{Full Feedback} serves as an oracle upper bound. The best-performing method for each column is highlighted in bold.}
% \label{tab:sprout_BERT_main}
% \resizebox{\linewidth}{!}{%
% \begin{tabular}{lccccccccccc}
% \toprule
% \textbf{Method} & $\lambda=0$ & $\lambda=100$ & $\lambda=200$ & $\lambda=300$ & $\lambda=400$ & $\lambda=500$ & $\lambda=600$ & $\lambda=700$ & $\lambda=800$ & $\lambda=900$ & $\lambda=1000$ \\
% \midrule
% Full-Feedback & \textit{85.99}$_{\pm0.17}$ & \textit{79.34}$_{\pm0.13}$ & \textit{75.55}$_{\pm0.29}$ & \textit{72.16}$_{\pm0.24}$ & \textit{69.47}$_{\pm0.23}$ & \textit{66.97}$_{\pm0.37}$ & \textit{64.79}$_{\pm0.31}$ & \textit{63.18}$_{\pm0.21}$ & \textit{61.43}$_{\pm0.27}$ & \textit{59.98}$_{\pm0.30}$ & \textit{58.83}$_{\pm0.24}$ \\ \hdashline
% Baseline & 66.88$_{\pm1.55}$ & 64.62$_{\pm0.66}$ & 61.85$_{\pm0.78}$ & 59.15$_{\pm0.84}$ & 56.73$_{\pm0.67}$ & 54.08$_{\pm0.57}$ & 51.17$_{\pm0.54}$ & 48.79$_{\pm0.67}$ & 45.74$_{\pm1.21}$ & 42.58$_{\pm1.36}$ & 38.97$_{\pm2.65}$ \\
% R\&C & 83.34$_{\pm0.32}$ & 76.45$_{\pm0.53}$ & 72.40$_{\pm0.56}$ & 68.59$_{\pm0.69}$ & 65.34$_{\pm0.66}$ & 63.19$_{\pm0.83}$ & 60.98$_{\pm0.81}$ & 59.01$_{\pm0.86}$ & 57.21$_{\pm0.99}$ & 55.97$_{\pm1.11}$ & 54.33$_{\pm1.37}$ \\
% CARROT (kNN) & 84.52$_{\pm0.35}$ & 76.55$_{\pm0.37}$ & 71.34$_{\pm0.23}$ & 67.82$_{\pm0.43}$ & 65.38$_{\pm0.40}$ & 63.38$_{\pm0.43}$ & 61.77$_{\pm0.43}$ & 60.39$_{\pm0.38}$ & 59.07$_{\pm0.37}$ & 57.99$_{\pm0.35}$ & 56.99$_{\pm0.34}$ \\
% Causal-CARROT (EmbedNet) & 83.46$_{\pm0.39}$ & 68.73$_{\pm0.74}$ & 62.44$_{\pm2.40}$ & 56.72$_{\pm1.68}$ & 54.37$_{\pm2.12}$ & 48.91$_{\pm2.41}$ & 46.35$_{\pm1.53}$ & 43.66$_{\pm3.04}$ & 41.81$_{\pm2.34}$ & 37.42$_{\pm5.86}$ & 34.69$_{\pm4.44}$ \\
% CF-Regression & 83.54$_{\pm0.25}$ & 76.65$_{\pm0.46}$ & 72.42$_{\pm0.53}$ & 68.95$_{\pm0.60}$ & 65.81$_{\pm0.67}$ & 63.58$_{\pm0.67}$ & 61.37$_{\pm0.64}$ & 59.52$_{\pm0.54}$ & 57.72$_{\pm0.72}$ & 56.31$_{\pm0.66}$ & 54.56$_{\pm0.71}$ \\
% Classification & 84.20$_{\pm0.24}$ & \textbf{77.58}$_{\pm0.34}$ & 73.36$_{\pm0.49}$ & 69.84$_{\pm0.48}$ & 66.49$_{\pm0.63}$ & 64.03$_{\pm1.02}$ & 61.84$_{\pm1.13}$ & 59.68$_{\pm1.48}$ & 58.09$_{\pm1.22}$ & 56.52$_{\pm1.63}$ & 54.85$_{\pm1.76}$ \\
% RM-Softmax & \textbf{84.97}$_{\pm0.39}$ & 77.53$_{\pm0.81}$ & \textbf{73.89}$_{\pm0.00}$ & \textbf{70.47}$_{\pm0.00}$ & \textbf{67.38}$_{\pm0.47}$ & \textbf{65.51}$_{\pm0.60}$ & \textbf{64.03}$_{\pm0.39}$ & \textbf{62.04}$_{\pm1.17}$ & \textbf{60.32}$_{\pm1.31}$ & \textbf{58.85}$_{\pm1.13}$ & \textbf{56.95}$_{\pm0.55}$ \\
% RM-Interval & \textbf{84.97}$_{\pm0.39}$ & 77.60$_{\pm0.62}$ & \textbf{73.89}$_{\pm0.00}$ & 69.92$_{\pm0.57}$ & \textbf{67.38}$_{\pm0.47}$ & 65.20$_{\pm0.67}$ & \textbf{64.03}$_{\pm0.39}$ & 61.54$_{\pm1.31}$ & \textbf{60.32}$_{\pm1.31}$ & 58.58$_{\pm1.15}$ & \textbf{56.95}$_{\pm0.55}$ \\
% \bottomrule
% \end{tabular}
% }
% \end{table}

\begin{figure}[h]
\centering
\includegraphics[width=0.94\linewidth]{images/cost_vs_accuracy_sprout_BERT.pdf}
\caption{Accuracy–cost trade-off curve for SPROUT with BERT embeddings.}
\label{fig:BERT_sprout}
\end{figure}


\begin{table}[h]
\centering
\caption{Utility on \textbf{SPROUT} with \texttt{BERT} embeddings. \textit{Full Feedback} serves as an oracle upper bound. The best-performing method for each column is highlighted in bold.}
\label{tab:sprout_BERT_main}
\resizebox{\linewidth}{!}{%
\begin{tabular}{lcccccc}
\toprule
\textbf{Method} & $\lambda=0$ & $\lambda=100$ & $\lambda=200$ & $\lambda=300$ & $\lambda=400$ & $\lambda=500$ \\
\midrule
Full-Feedback & \textit{85.99}$_{\pm0.17}$ & \textit{79.34}$_{\pm0.13}$ & \textit{75.55}$_{\pm0.29}$ & \textit{72.16}$_{\pm0.24}$ & \textit{69.47}$_{\pm0.23}$ & \textit{66.97}$_{\pm0.37}$ \\ \hdashline
Baseline & 66.88$_{\pm1.55}$ & 64.62$_{\pm0.66}$ & 61.85$_{\pm0.78}$ & 59.15$_{\pm0.84}$ & 56.73$_{\pm0.67}$ & 54.08$_{\pm0.57}$ \\
R\&C & 83.34$_{\pm0.32}$ & 76.45$_{\pm0.53}$ & 72.40$_{\pm0.56}$ & 68.59$_{\pm0.69}$ & 65.34$_{\pm0.66}$ & 63.19$_{\pm0.83}$ \\
Causal-CARROT (kNN) & 84.52$_{\pm0.35}$ & 76.55$_{\pm0.37}$ & 71.34$_{\pm0.23}$ & 67.82$_{\pm0.43}$ & 65.38$_{\pm0.40}$ & 63.38$_{\pm0.43}$ \\
Causal-CARROT (EmbedNet) & 83.46$_{\pm0.39}$ & 68.73$_{\pm0.74}$ & 62.44$_{\pm2.40}$ & 56.72$_{\pm1.68}$ & 54.37$_{\pm2.12}$ & 48.91$_{\pm2.41}$ \\
CF-Regression & 83.54$_{\pm0.25}$ & 76.65$_{\pm0.46}$ & 72.42$_{\pm0.53}$ & 68.95$_{\pm0.60}$ & 65.81$_{\pm0.67}$ & 63.58$_{\pm0.67}$ \\
RM-Classification & 84.20$_{\pm0.24}$ & \textbf{77.58}$_{\pm0.34}$ & 73.36$_{\pm0.49}$ & 69.84$_{\pm0.48}$ & 66.49$_{\pm0.63}$ & 64.03$_{\pm1.02}$ \\
RM-Softmax & \textbf{84.97}$_{\pm0.39}$ & 77.53$_{\pm0.81}$ & \textbf{73.89}$_{\pm0.00}$ & \textbf{70.47}$_{\pm0.00}$ & \textbf{67.38}$_{\pm0.47}$ & \textbf{65.51}$_{\pm0.60}$ \\ \hdashline
RM-Interval & \textbf{84.97}$_{\pm0.39}$ & 77.60$_{\pm0.62}$ & \textbf{73.89}$_{\pm0.00}$ & 69.92$_{\pm0.57}$ & \textbf{67.38}$_{\pm0.47}$ & 65.20$_{\pm0.67}$ \\
\midrule
\textbf{Method} & $\lambda=600$ & $\lambda=700$ & $\lambda=800$ & $\lambda=900$ & $\lambda=1000$ & \\
\midrule
Full-Feedback & \textit{64.79}$_{\pm0.31}$ & \textit{63.18}$_{\pm0.21}$ & \textit{61.43}$_{\pm0.27}$ & \textit{59.98}$_{\pm0.30}$ & \textit{58.83}$_{\pm0.24}$ & \\ \hdashline
Baseline & 51.17$_{\pm0.54}$ & 48.79$_{\pm0.67}$ & 45.74$_{\pm1.21}$ & 42.58$_{\pm1.36}$ & 38.97$_{\pm2.65}$ & \\
R\&C & 60.98$_{\pm0.81}$ & 59.01$_{\pm0.86}$ & 57.21$_{\pm0.99}$ & 55.97$_{\pm1.11}$ & 54.33$_{\pm1.37}$ & \\
Causal-CARROT (kNN) & 61.77$_{\pm0.43}$ & 60.39$_{\pm0.38}$ & 59.07$_{\pm0.37}$ & 57.99$_{\pm0.35}$ & 56.99$_{\pm0.34}$ & \\
Causal-CARROT (EmbedNet) & 46.35$_{\pm1.53}$ & 43.66$_{\pm3.04}$ & 41.81$_{\pm2.34}$ & 37.42$_{\pm5.86}$ & 34.69$_{\pm4.44}$ & \\
CF-Regression & 61.37$_{\pm0.64}$ & 59.52$_{\pm0.54}$ & 57.72$_{\pm0.72}$ & 56.31$_{\pm0.66}$ & 54.56$_{\pm0.71}$ & \\
RM-Classification & 61.84$_{\pm1.13}$ & 59.68$_{\pm1.48}$ & 58.09$_{\pm1.22}$ & 56.52$_{\pm1.63}$ & 54.85$_{\pm1.76}$ & \\
RM-Softmax & \textbf{64.03}$_{\pm0.39}$ & \textbf{62.04}$_{\pm1.17}$ & \textbf{60.32}$_{\pm1.31}$ & \textbf{58.85}$_{\pm1.13}$ & \textbf{56.95}$_{\pm0.55}$ & \\ \hdashline
RM-Interval & \textbf{64.03}$_{\pm0.39}$ & 61.54$_{\pm1.31}$ & \textbf{60.32}$_{\pm1.31}$ & 58.58$_{\pm1.15}$ & \textbf{56.95}$_{\pm0.55}$ & \\
\bottomrule
\end{tabular}
}
\end{table}







Our RM-based approaches consistently deliver the strongest performance overall, with \textbf{RM-Softmax} and \textbf{RM-Interval} standing out for both high utility and low variance. Notably, \textbf{RM-Interval} generalizes remarkably well to unseen budget levels (i.e., odd  $\lambda$ values), many times even outperforming models trained specifically on those points. These results underscore the effectiveness of our regret-minimization framework in both fixed and variable cost settings.

The standard \textbf{Baseline} method, which reflects the common decoupled approach used in prior work and ignores treatment selection bias, performs the worst across most values of $\lambda$, underscoring the importance of accounting for treatment bias in observational data. Notably,  the simple \textbf{R\&C} method and the \textbf{Causal-CARROT} variants (which incorporate causal corrections) achieve substantial improvements over it, validating our claim that bias-aware routing significantly improves performance.

The performance gap between \textbf{CF-Regression} and our RM-based methods demonstrates the benefit of an integrated end-to-end approach. Whereas \textbf{CF-Regression} focuses on approximating the counterfactual utility function and then selecting the best model based on predicted outcomes, our methods directly minimize regret, leading to superior and more stable results.
Comparing our two surrogate formulations, \textbf{RM-Softmax} generally outperforms \textbf{RM-Classification} in both utility and variance, indicating the advantage of optimizing a differentiable surrogate objective. Finally, among the \textbf{Causal-CARROT} variants, the {kNN} version consistently outperforms the EmbedNet variant, suggesting that non-parametric estimators may offer greater robustness in this setting.

%Across all settings, our regret-based, classification-based and interpolation-based methods outperform existing baselines by a margin. The regret-optimized router consistently achieves the highest utility, demonstrating strong performance across both datasets and embedding models. Its performance is closely matched by InterpolationNet, which additionally supports continuous generalization to unseen budget levels (\(\lambda \in \{100, 300, \dots, 900\}\)). Our classification-based approach also performs competitively, consistently outperforming both CARROT variants and the R\&C estimator.

%Among benchmarks, CARROT (kNN) is the strongest, consistently outperforming its learned embedding variant. The R\&C approach yields notable improvements over the baselines, underscoring the importance of correcting for treatment assignment bias. Furthermore, the sequential regression-based model performs reasonably well, our end-to-end approaches achieve superior results, highlighting the benefits of directly optimizing for decision quality rather than relying on intermediate predictions.


\section{Conclusion}

We propose a causal end-to-end framework for routing queries to LLMs under observational data. Our approach introduces a regret-minimizing objective grounded in counterfactual estimation, enabling principled policy learning that accounts for treatment selection bias without requiring full-feedback data. Unlike prior approaches that rely on decoupled prediction of accuracy and cost, where errors can compound, our method directly optimizes the decision objective. To support heterogeneous user preferences, we develop an interval-conditioned routing architecture that generalizes across a continuum of cost-sensitivity parameters. Theoretical analysis provides guarantees on interpolation sufficiency and regret bounds, while empirical evaluations on public routing benchmarks demonstrate that our methods consistently outperform strong baselines, including recent routing algorithms, across multiple embedding models. Future work includes extending the framework to accommodate additional user-defined metrics or hard constraints that cannot be readily incorporated as soft penalties in the objective. Another promising direction is to explore online or adaptive routing in dynamic environments, as well as extending causal regret minimization to multi-turn settings.

% \newpage 
\nocite{*}
\bibliographystyle{plainnat}
\bibliography{references}







\newpage
\appendix
\title{app}

\setcounter{proposition}{0}
\setcounter{corollary}{0}

% Redefine numbering if needed (optional)
\renewcommand{\theproposition}{\arabic{proposition}}
\renewcommand{\thecorollary}{\arabic{corollary}}





\section{Proofs of Propositions}

\begin{proposition}
Suppose the estimated utility function $\hat{Y}_x : \mathcal{T} \to \mathbb{R}$ is $L$-Lipschitz continuous over the probability simplex with respect to the $\ell_1$ norm, as in Definition~\ref{def:lip}. Then, for a policy $f : \mathcal{X} \to \Delta^{|\mathcal{T}|}$ that outputs a distribution $f(x)$ over $\mathcal{T}$, the regret can be upper bounded by:
\begin{align}
    \text{Regret}(f) \leq L \cdot \frac{1}{n} \sum_{i=1}^n \sqrt{2 \cdot \text{CE}(t_i^*, f(x_i))},
\end{align}
where $t_i^* := \arg\max_{t \in \mathcal{T}} \hat{Y}_{x_i}(t)$ is the optimal treatment for input $x_i$, and $\text{CE}(t_i^*, f(x_i)) := -\log f(x_i)_{t_i^*}$ denotes the cross-entropy loss.
\end{proposition}

\begin{proof}
Let $e_{t_i^*} \in \Delta^{|\mathcal{T}|}$ denote the one-hot distribution over the optimal treatment $t_i^*$. By the definition of regret:
\begin{align}
\text{Regret}(f) &= \frac{1}{n} \sum_{i=1}^n \left[ \hat{Y}_{x_i}(e_{t_i^*}) - \hat{Y}_{x_i}(f(x_i)) \right].
\end{align}
\noindent Using the $L$-Lipschitz continuity of $\hat{Y}_{x_i}(\cdot)$ under the $\ell_1$ norm:
\begin{align}
\left| \hat{Y}_{x_i}(e_{t_i^*}) - \hat{Y}_{x_i}(f(x_i)) \right| &\leq L \cdot \|e_{t_i^*} - f(x_i)\|_1.
\end{align}
\noindent Applying Pinsker’s inequality:
\begin{align}
\|e_{t_i^*} - f(x_i)\|_1 &\leq \sqrt{2 \cdot \text{KL}(e_{t_i^*} \,\|\, f(x_i))} = \sqrt{2 \cdot \text{CE}(t_i^*, f(x_i))},
\end{align}
where the equality follows because KL divergence from a one-hot distribution to a probability vector reduces to cross-entropy. Combining the above:
\begin{align}
\text{Regret}(f) &\leq L \cdot \frac{1}{n} \sum_{i=1}^n \sqrt{2 \cdot \text{CE}(t_i^*, f(x_i))},
\end{align}
which completes the proof.
\end{proof}



\begin{proposition}
Let $f: \mathcal{X} \to \mathbb{R}^{|\mathcal{T}|}$ be a neural network whose output is passed through a softmax layer with fixed temperature $\tau > 0$, and define $t^*_i := \arg\max_{t \in \mathcal{T}} \hat{Y}_{x_i}(t)$. Then, optimizing  the following objective using gradient descent 
\begin{align}
    \min_f \frac{1}{n} \sum_{i=1}^n \left( \hat{Y}_{x_i}(t^*_i) - \sum_{t=1}^{|\mathcal{T}|} \hat{Y}_{x_i}(t) \cdot \mathrm{softmax}(f(x_i))_t \right)
\end{align}
leads the model $f$ to place all probability mass on the optimal treatment $t^*_i$. That is, at convergence,
\begin{align}
    \mathrm{softmax}(f(x_i))_t \to 
    \begin{cases}
        1 & \text{if } t = t^*_i, \\
        0 & \text{otherwise}.
    \end{cases}
\end{align}
\end{proposition}

\begin{proof}
Let $\hat{Y}_{x_i} \in \mathbb{R}^{|\mathcal{T}|}$ denote the vector of estimated potential outcomes for input $x_i$, and let $f(x_i) \in \mathbb{R}^{|\mathcal{T}|}$ be the output of the neural network before the softmax layer. The objective for a single instance $x_i$ can be written as minimizing the regret surrogate:
\begin{align}
    \hat{Y}_{x_i}(t^*_i) - \sum_{t=1}^{|\mathcal{T}|} \hat{Y}_{x_i}(t) \cdot \mathrm{softmax}(f(x_i))_t.
\end{align}
This is equivalent to maximizing the inner product:
\begin{align}
    \langle \hat{Y}_{x_i}, \mathrm{softmax}(f(x_i)) \rangle.
\end{align}

Let us denote $p := \mathrm{softmax}(f(x_i)) \in \Delta^{|\mathcal{T}| - 1}$, the probability simplex. We now show that the inner product $\langle \hat{Y}_{x_i}, p \rangle$ increases at each gradient step. Since $p = \mathrm{softmax}(f(x_i))$, we can compute the gradient of the objective with respect to $f(x_i)$ as:
\begin{align}
    \nabla_{f(x_i)} \langle \hat{Y}_{x_i}, \mathrm{softmax}(f(x_i)) \rangle
    = J_{\mathrm{softmax}}(f(x_i))^\top \hat{Y}_{x_i},
\end{align}
where $J_{\mathrm{softmax}}(f(x_i))$ is the Jacobian of the softmax function, given by:
\begin{align}
    J_{\mathrm{softmax}}(f(x_i))_{t,s} = \frac{\partial \, \mathrm{softmax}(f(x_i))_t}{\partial f(x_i)_s}
    = \mathrm{softmax}(f(x_i))_t \left(\delta_{t,s} - \mathrm{softmax}(f(x_i))_s \right).
\end{align}

This gradient direction corresponds to increasing the logit value of actions with higher $\hat{Y}_{x_i}(t)$ and decreasing those with lower values, pushing the softmax distribution toward the mode of $\hat{Y}_{x_i}$. In other words, the gradient ascent step increases the inner product at each iteration $k$:
\begin{align}
    \langle \hat{Y}_{x_i}, \mathrm{softmax}(f(x_i)) \rangle^{(k+1)} 
    > \langle \hat{Y}_{x_i}, \mathrm{softmax}(f(x_i)) \rangle^{(k)}.
\end{align}

Since $\hat{Y}_{x_i}$ is fixed and the softmax is smooth and bounded, this sequence is monotonically increasing and converges to the maximum possible value:
\begin{align}
    \langle \hat{Y}_{x_i}, \mathrm{softmax}(f(x_i)) \rangle \to \max_t \hat{Y}_{x_i}(t) = \hat{Y}_{x_i}(t^*_i),
\end{align}
which implies:
\begin{align}
    \mathrm{softmax}(f(x_i))_t \to 
    \begin{cases}
        1 & \text{if } t = t^*_i, \\
        0 & \text{otherwise}.
    \end{cases}
\end{align}

Thus, the regret surrogate converges to zero:
\begin{align}
    \hat{Y}_{x_i}(t^*_i) - \langle \hat{Y}_{x_i}, \mathrm{softmax}(f(x_i)) \rangle \to 0,
\end{align}
and the learned policy selects the treatment maximizing the estimated outcome.
\end{proof}

\begin{proposition}[Piecewise Constant Optimal Policy]
Fix a query $x \in \mathcal{X}$ and assume the estimated utility function is affine in $\lambda$, i.e., $\hat{Y}_x^\lambda(t) = a_x(t) - \lambda \cdot c_x(t)$ for all $t \in \mathcal{T}$. Then the optimal treatment
\[
t^*(\lambda) := \arg\max_{t \in \mathcal{T}} \hat{Y}_x^\lambda(t)
\]
is piecewise constant in $\lambda$. That is, the budget space $\mathbb{R}_{\geq 0}$ can be partitioned into intervals over which the optimal treatment remains fixed.
\end{proposition}

\begin{proof}
For fixed $x$, each $\hat{Y}_x^\lambda(t)$ is an affine function of $\lambda$. The pointwise maximum of a finite collection of affine functions is piecewise affine, and the argmax corresponds to the highest line at each $\lambda$. Since each pair of lines can intersect at most once, the number of intervals over which a single treatment is optimal is bounded by $|\mathcal{T}| - 1$. Therefore, $t^*(\lambda)$ changes only at these intersection points and remains constant within each interval.
\end{proof}

\begin{proposition}[Affine Closure of Utility Function]
Let $\lambda_j < \lambda_{j+1}$ be two adjacent budget values and let $\lambda \in [\lambda_j, \lambda_{j+1}]$. Suppose the utility function is affine in $\lambda$:
\[
\hat{Y}_x^\lambda(t) = a_x(t) - \lambda \cdot c_x(t).
\]
Then for all $t \in \mathcal{T}$, the utility at $\lambda$ is a convex combination of utilities at the endpoints:
\[
\hat{Y}_x^\lambda(t) = \alpha \cdot \hat{Y}_x^{\lambda_j}(t) + (1 - \alpha) \cdot \hat{Y}_x^{\lambda_{j+1}}(t),
\quad \text{where } \alpha := \frac{\lambda_{j+1} - \lambda}{\lambda_{j+1} - \lambda_j}.
\]
\end{proposition}

\begin{proof}
We expand each term:
\begin{align*}
\hat{Y}_x^{\lambda_j}(t) &= a_x(t) - \lambda_j \cdot c_x(t), \\
\hat{Y}_x^{\lambda_{j+1}}(t) &= a_x(t) - \lambda_{j+1} \cdot c_x(t).
\end{align*}
Then:
\begin{align*}
\alpha \cdot \hat{Y}_x^{\lambda_j}(t) + (1 - \alpha) \cdot \hat{Y}_x^{\lambda_{j+1}}(t)
&= \alpha \cdot (a_x(t) - \lambda_j c_x(t)) + (1 - \alpha) \cdot (a_x(t) - \lambda_{j+1} c_x(t)) \\
&= a_x(t) - [\alpha \lambda_j + (1 - \alpha) \lambda_{j+1}] \cdot c_x(t) \\
&= a_x(t) - \lambda \cdot c_x(t) = \hat{Y}_x^\lambda(t),
\end{align*}
since:
\[
\alpha \lambda_j + (1 - \alpha) \lambda_{j+1} = \lambda.
\]
\end{proof}

\begin{corollary}[Sufficiency of Two Models per Interval]
Under the affine assumption, the utility $\hat{Y}_x^\lambda(t)$ for any $\lambda \in [\lambda_j, \lambda_{j+1}]$ can be exactly reconstructed using only the endpoints $\hat{Y}_x^{\lambda_j}(t)$ and $\hat{Y}_x^{\lambda_{j+1}}(t)$. Thus, it is sufficient to use only the two corresponding models $f_{\lambda_j}$ and $f_{\lambda_{j+1}}$ for interpolation within the interval.
\end{corollary}

\begin{proof}
    This follows immediately from the statement of Proposition 4.
\end{proof}

\begin{proposition}[Expressivity of Additive Two-Model joint Architecture]
Let $\lambda \in [\lambda_j, \lambda_{j+1}]$, and suppose that for each $t \in \mathcal{T}$ the utility function satisfies $\hat{Y}_x^\lambda(t) = a_x(t) - \lambda \cdot c_x(t)$. Then the optimal treatment $t^*(\lambda) := \arg\max_{t} \hat{Y}_x^\lambda(t)$ can be exactly represented by a softmax policy over a function of the form:
\[
f(x, \lambda) = \texttt{Linear} \left( [f_{\lambda_j}(x), f_{\lambda_{j+1}}(x)] + g(\lambda) \right),
\]
where $g(\lambda)$ is any differentiable embedding of $\lambda$, and $f_{\lambda_j}, f_{\lambda_{j+1}}$ are accurate predictors trained at endpoints $\lambda_j$ and $\lambda_{j+1}$.
\end{proposition}


\begin{proof}
From Proposition~\ref{prop:affine-closure}, the utility $\hat{Y}_x^\lambda(t)$ is a convex combination of $\hat{Y}_x^{\lambda_j}(t)$ and $\hat{Y}_x^{\lambda_{j+1}}(t)$. If the network $f(x, \lambda)$ linearly combines the outputs of $f_{\lambda_j}(x)$ and $f_{\lambda_{j+1}}(x)$, then its scores can match $\hat{Y}_x^\lambda(t)$ up to a scalar transformation. Applying softmax preserves the argmax.

Including $g(\lambda)$ allows the architecture to learn any additional monotonic reweighting of the interpolation, ensuring the output scores can be shaped to approximate the true utility surface exactly. Thus, the architecture can represent the optimal policy within each interval.
\end{proof}


\section{Additional Results}\label{app:single_lambda_results}

In this section, we present the additional plots for the rest of the datasets as well as the exact values of utility for value $\lambda$. We begin by presenting the rest of the figures.

\subsection{Additional Figures}
\begin{figure}[H]
\centering
\includegraphics[width=1\linewidth]{images/cost_vs_accuracy_routerbench_BERT.pdf}
\caption{Accuracy–cost trade-off curve for RouterBench with BERT embeddings.}
\end{figure}

\begin{figure}[H]
\centering
\includegraphics[width=1\linewidth]{images/cost_vs_accuracy_routerbench_llama-3.2-1B.pdf}
\caption{Accuracy–cost trade-off curve for RouterBench with LLaMa-3.2-1B embeddings.}
\label{fig:single_lambda_plot}
\end{figure}

\begin{figure}[H]
\centering
\includegraphics[width=1\linewidth]{images/cost_vs_accuracy_sprout_BERT.pdf}
\caption{Accuracy–cost trade-off curve for SPROUT with BERT embeddings.}
\end{figure}


\begin{figure}[H]
\centering
\includegraphics[width=1\linewidth]{images/cost_vs_accuracy_sprout_llama-3.2-1B.pdf}
\caption{Accuracy–cost trade-off curve for SPROUT with LLaMa-3.2-1B embeddings.}
\end{figure}



\subsection{Additional Tables}

\begin{table}[H]
\centering
\caption{Utility for \textbf{RouterBench} with \texttt{BERT} embeddings. \textit{Full Feedback} serves as an oracle upper bound. The best-performing method for each column is highlighted in bold.}
\label{tab:routerbench_bert}
\resizebox{\linewidth}{!}{%
\begin{tabular}{lccccccccccc}
\toprule
\textbf{Method} & $\lambda=0$ & $\lambda=100$ & $\lambda=200$ & $\lambda=300$ & $\lambda=400$ & $\lambda=500$ & $\lambda=600$ & $\lambda=700$ & $\lambda=800$ & $\lambda=900$ & $\lambda=1000$ \\
\midrule
Full-Feedback & \textit{78.07}$_{\pm0.14}$ & \textit{68.07}$_{\pm0.24}$ & \textit{64.12}$_{\pm0.18}$ & \textit{61.96}$_{\pm0.13}$ & \textit{60.31}$_{\pm0.13}$ & \textit{58.92}$_{\pm0.18}$ & \textit{57.62}$_{\pm0.21}$ & \textit{56.26}$_{\pm0.30}$ & \textit{55.00}$_{\pm0.24}$ & \textit{53.73}$_{\pm0.31}$ & \textit{52.42}$_{\pm0.18}$ \\ \hdashline
Baseline & 66.46$_{\pm0.66}$ & 61.60$_{\pm0.64}$ & 59.75$_{\pm0.46}$ & 57.41$_{\pm0.49}$ & 55.37$_{\pm0.51}$ & 53.85$_{\pm0.96}$ & 51.67$_{\pm1.04}$ & 50.55$_{\pm1.07}$ & 48.72$_{\pm1.28}$ & 47.87$_{\pm0.76}$ & 46.19$_{\pm0.46}$ \\
R\&C & 75.08$_{\pm0.49}$ & 65.58$_{\pm0.28}$ & 61.73$_{\pm0.53}$ & 59.63$_{\pm0.53}$ & 58.23$_{\pm0.58}$ & 56.79$_{\pm0.63}$ & 55.47$_{\pm0.64}$ & 54.34$_{\pm0.46}$ & 53.08$_{\pm0.58}$ & 51.65$_{\pm0.59}$ & 50.18$_{\pm0.53}$ \\
Causal-CARROT (kNN) & 75.12$_{\pm0.64}$ & 65.12$_{\pm0.51}$ & 62.21$_{\pm0.52}$ & 60.56$_{\pm0.49}$ & 58.94$_{\pm0.41}$ & 57.39$_{\pm0.42}$ & 55.86$_{\pm0.35}$ & 54.41$_{\pm0.42}$ & 52.95$_{\pm0.46}$ & 51.47$_{\pm0.49}$ & 50.08$_{\pm0.60}$ \\
Causal-CARROT (EmbedNet) & 75.07$_{\pm0.50}$ & 56.87$_{\pm1.39}$ & 47.71$_{\pm4.76}$ & 44.66$_{\pm1.95}$ & 42.06$_{\pm2.31}$ & 38.55$_{\pm3.30}$ & 35.22$_{\pm2.27}$ & 32.25$_{\pm2.40}$ & 26.66$_{\pm6.05}$ & 28.89$_{\pm3.53}$ & 24.64$_{\pm5.52}$ \\
CF-Regression & 74.81$_{\pm0.55}$ & 65.56$_{\pm0.17}$ & 61.95$_{\pm0.44}$ & 59.71$_{\pm0.69}$ & 57.91$_{\pm0.65}$ & 56.36$_{\pm0.73}$ & 54.87$_{\pm0.55}$ & 53.44$_{\pm0.63}$ & 52.19$_{\pm0.48}$ & 50.78$_{\pm0.51}$ & 49.43$_{\pm0.39}$ \\
RM-Classification & 76.30$_{\pm0.61}$ & \textbf{66.43}$_{\pm0.28}$ & 62.65$_{\pm0.64}$ & 60.85$_{\pm0.54}$ & \textbf{59.65}$_{\pm0.51}$ & \textbf{58.07}$_{\pm0.62}$ & \textbf{56.80}$_{\pm0.47}$ & \textbf{55.29}$_{\pm0.58}$ & \textbf{54.07}$_{\pm0.62}$ & \textbf{52.45}$_{\pm0.72}$ & \textbf{51.14}$_{\pm0.61}$ \\
RM-Softmax & \textbf{77.82}$_{\pm0.03}$ & 65.51$_{\pm0.51}$ & \textbf{63.30}$_{\pm0.26}$ & 60.75$_{\pm0.61}$ & 58.58$_{\pm0.59}$ & 56.48$_{\pm0.52}$ & 54.53$_{\pm0.57}$ & 52.79$_{\pm0.74}$ & 51.10$_{\pm0.98}$ & 49.22$_{\pm1.14}$ & 47.64$_{\pm1.40}$ \\ \hdashline
RM-Interval & \textbf{77.82}$_{\pm0.03}$ & 65.51$_{\pm0.30}$ & \textbf{63.30}$_{\pm0.26}$ & \textbf{61.01}$_{\pm0.51}$ & 58.58$_{\pm0.59}$ & 56.92$_{\pm0.31}$ & 54.53$_{\pm0.57}$ & 53.11$_{\pm0.77}$ & 51.10$_{\pm0.98}$ & 49.60$_{\pm1.21}$ & 47.64$_{\pm1.40}$ \\
\bottomrule
\end{tabular}
}
\end{table}

\begin{table}[H]
\centering
\caption{Utility for \textbf{RouterBench} with \texttt{Llama-3.2-1B} embeddings. %\textit{Full Feedback} serves as an oracle upper bound. 
The best-performing method for each column is highlighted in bold.}
\label{tab:routerbench_llama-3.2-1B}
\resizebox{\linewidth}{!}{%
\begin{tabular}{lccccccccccc}
\toprule
\textbf{Method} & $\lambda=0$ & $\lambda=100$ & $\lambda=200$ & $\lambda=300$ & $\lambda=400$ & $\lambda=500$ & $\lambda=600$ & $\lambda=700$ & $\lambda=800$ & $\lambda=900$ & $\lambda=1000$ \\
\midrule
Full-Feedback & \textit{77.06}$_{\pm0.30}$ & \textit{67.37}$_{\pm0.35}$ & \textit{63.52}$_{\pm0.35}$ & \textit{60.94}$_{\pm0.27}$ & \textit{59.70}$_{\pm0.24}$ & \textit{58.20}$_{\pm0.19}$ & \textit{56.88}$_{\pm0.33}$ & \textit{55.49}$_{\pm0.38}$ & \textit{54.03}$_{\pm0.14}$ & \textit{52.86}$_{\pm0.25}$ & \textit{51.47}$_{\pm0.28}$ \\ \hdashline
Baseline & 50.24$_{\pm1.28}$ & 54.56$_{\pm0.96}$ & 54.79$_{\pm1.11}$ & 53.54$_{\pm0.63}$ & 52.00$_{\pm0.57}$ & 50.63$_{\pm0.68}$ & 49.43$_{\pm0.66}$ & 47.78$_{\pm1.22}$ & 46.19$_{\pm0.61}$ & 44.96$_{\pm0.66}$ & 43.10$_{\pm0.84}$ \\
R\&C & 69.30$_{\pm0.63}$ & 61.62$_{\pm0.70}$ & 58.04$_{\pm0.44}$ & 56.20$_{\pm0.52}$ & 54.33$_{\pm0.58}$ & 53.02$_{\pm0.62}$ & 51.72$_{\pm0.57}$ & 50.30$_{\pm0.50}$ & 48.94$_{\pm0.48}$ & 48.03$_{\pm0.56}$ & 46.73$_{\pm0.56}$ \\
Causal-CARROT (kNN) & 75.04$_{\pm0.47}$ & 64.72$_{\pm0.47}$ & 61.86$_{\pm0.53}$ & \textbf{60.23}$_{\pm0.54}$ & \textbf{58.71}$_{\pm0.49}$ & \textbf{57.18}$_{\pm0.48}$ & \textbf{55.58}$_{\pm0.50}$ & \textbf{54.04}$_{\pm0.49}$ & \textbf{52.50}$_{\pm0.51}$ & \textbf{51.00}$_{\pm0.59}$ & \textbf{49.54}$_{\pm0.54}$ \\
Causal-CARROT (EmbedNet) & 69.30$_{\pm0.38}$ & 44.62$_{\pm1.38}$ & 38.94$_{\pm2.63}$ & 33.82$_{\pm2.07}$ & 28.60$_{\pm4.29}$ & 24.04$_{\pm3.76}$ & 20.47$_{\pm6.44}$ & 15.11$_{\pm3.78}$ & 12.32$_{\pm5.07}$ & 9.64$_{\pm6.26}$ & 13.48$_{\pm1.45}$ \\
CF-Regression & 70.02$_{\pm0.33}$ & 62.33$_{\pm0.78}$ & 59.05$_{\pm0.24}$ & 57.19$_{\pm0.42}$ & 55.58$_{\pm0.56}$ & 54.08$_{\pm0.43}$ & 53.03$_{\pm0.64}$ & 51.97$_{\pm0.64}$ & 50.50$_{\pm0.56}$ & 49.10$_{\pm0.45}$ & 47.68$_{\pm0.44}$ \\
RM-Classification & 72.76$_{\pm0.50}$ & 64.22$_{\pm0.65}$ & 60.68$_{\pm0.56}$ & 58.54$_{\pm0.46}$ & 56.64$_{\pm0.72}$ & 55.44$_{\pm0.33}$ & 53.98$_{\pm0.46}$ & 52.71$_{\pm0.44}$ & 51.14$_{\pm0.59}$ & 50.06$_{\pm0.55}$ & 48.63$_{\pm0.36}$ \\
RM-Softmax & \textbf{77.58}$_{\pm0.49}$ & \textbf{64.49}$_{\pm0.50}$ & \textbf{61.97}$_{\pm0.52}$ & 60.17$_{\pm0.61}$ & 58.20$_{\pm0.46}$ & 56.21$_{\pm0.40}$ & 54.34$_{\pm0.51}$ & 52.86$_{\pm0.79}$ & 50.96$_{\pm0.92}$ & 49.70$_{\pm1.19}$ & 47.69$_{\pm1.65}$ \\ \hdashline
RM-Interval & \textbf{77.58}$_{\pm0.49}$ & 63.70$_{\pm1.47}$ & \textbf{61.97}$_{\pm0.52}$ & 59.43$_{\pm1.56}$ & 58.20$_{\pm0.46}$ & 56.31$_{\pm0.47}$ & 54.34$_{\pm0.51}$ & 52.55$_{\pm0.91}$ & 50.96$_{\pm0.92}$ & 49.23$_{\pm1.81}$ & 47.69$_{\pm1.65}$ \\
\bottomrule
\end{tabular}
}
\end{table}


\begin{table}[H]
\centering
\caption{Utility for \textbf{SPOUT} with \texttt{BERT} embeddings. %\textit{Full Feedback} serves as an oracle upper bound. 
The best-performing method for each column is highlighted in bold.}
\label{tab:sprout_bert}
\resizebox{\linewidth}{!}{%
\begin{tabular}{lccccccccccc}
\toprule
\textbf{Method} & $\lambda=0$ & $\lambda=100$ & $\lambda=200$ & $\lambda=300$ & $\lambda=400$ & $\lambda=500$ & $\lambda=600$ & $\lambda=700$ & $\lambda=800$ & $\lambda=900$ & $\lambda=1000$ \\
\midrule
Full-Feedback & \textit{85.99}$_{\pm0.17}$ & \textit{79.34}$_{\pm0.13}$ & \textit{75.55}$_{\pm0.29}$ & \textit{72.16}$_{\pm0.24}$ & \textit{69.47}$_{\pm0.23}$ & \textit{66.97}$_{\pm0.37}$ & \textit{64.79}$_{\pm0.31}$ & \textit{63.18}$_{\pm0.21}$ & \textit{61.43}$_{\pm0.27}$ & \textit{59.98}$_{\pm0.30}$ & \textit{58.83}$_{\pm0.24}$ \\ \hdashline
Baseline & 66.88$_{\pm1.55}$ & 64.62$_{\pm0.66}$ & 61.85$_{\pm0.78}$ & 59.15$_{\pm0.84}$ & 56.73$_{\pm0.67}$ & 54.08$_{\pm0.57}$ & 51.17$_{\pm0.54}$ & 48.79$_{\pm0.67}$ & 45.74$_{\pm1.21}$ & 42.58$_{\pm1.36}$ & 38.97$_{\pm2.65}$ \\
R\&C & 83.34$_{\pm0.32}$ & 76.45$_{\pm0.53}$ & 72.40$_{\pm0.56}$ & 68.59$_{\pm0.69}$ & 65.34$_{\pm0.66}$ & 63.19$_{\pm0.83}$ & 60.98$_{\pm0.81}$ & 59.01$_{\pm0.86}$ & 57.21$_{\pm0.99}$ & 55.97$_{\pm1.11}$ & 54.33$_{\pm1.37}$ \\
Causal-CARROT (kNN) & 84.52$_{\pm0.35}$ & 76.55$_{\pm0.37}$ & 71.34$_{\pm0.23}$ & 67.82$_{\pm0.43}$ & 65.38$_{\pm0.40}$ & 63.38$_{\pm0.43}$ & 61.77$_{\pm0.43}$ & 60.39$_{\pm0.38}$ & 59.07$_{\pm0.37}$ & 57.99$_{\pm0.35}$ & 56.99$_{\pm0.34}$ \\
Causal-CARROT (EmbedNet) & 83.46$_{\pm0.39}$ & 68.73$_{\pm0.74}$ & 62.44$_{\pm2.40}$ & 56.72$_{\pm1.68}$ & 54.37$_{\pm2.12}$ & 48.91$_{\pm2.41}$ & 46.35$_{\pm1.53}$ & 43.66$_{\pm3.04}$ & 41.81$_{\pm2.34}$ & 37.42$_{\pm5.86}$ & 34.69$_{\pm4.44}$ \\
CF-Regression & 83.54$_{\pm0.25}$ & 76.65$_{\pm0.46}$ & 72.42$_{\pm0.53}$ & 68.95$_{\pm0.60}$ & 65.81$_{\pm0.67}$ & 63.58$_{\pm0.67}$ & 61.37$_{\pm0.64}$ & 59.52$_{\pm0.54}$ & 57.72$_{\pm0.72}$ & 56.31$_{\pm0.66}$ & 54.56$_{\pm0.71}$ \\
RM-Classification & 84.20$_{\pm0.24}$ & \textbf{77.58}$_{\pm0.34}$ & 73.36$_{\pm0.49}$ & 69.84$_{\pm0.48}$ & 66.49$_{\pm0.63}$ & 64.03$_{\pm1.02}$ & 61.84$_{\pm1.13}$ & 59.68$_{\pm1.48}$ & 58.09$_{\pm1.22}$ & 56.52$_{\pm1.63}$ & 54.85$_{\pm1.76}$ \\
RM-Softmax & \textbf{84.97}$_{\pm0.39}$ & 77.53$_{\pm0.81}$ & \textbf{73.89}$_{\pm0.00}$ & \textbf{70.47}$_{\pm0.00}$ & \textbf{67.38}$_{\pm0.47}$ & \textbf{65.51}$_{\pm0.60}$ & \textbf{64.03}$_{\pm0.39}$ & \textbf{62.04}$_{\pm1.17}$ & \textbf{60.32}$_{\pm1.31}$ & \textbf{58.85}$_{\pm1.13}$ & \textbf{56.95}$_{\pm0.55}$ \\ \hdashline
RM-Interval & \textbf{84.97}$_{\pm0.39}$ & 77.60$_{\pm0.62}$ & \textbf{73.89}$_{\pm0.00}$ & 69.92$_{\pm0.57}$ & \textbf{67.38}$_{\pm0.47}$ & 65.20$_{\pm0.67}$ & \textbf{64.03}$_{\pm0.39}$ & 61.54$_{\pm1.31}$ & \textbf{60.32}$_{\pm1.31}$ & 58.58$_{\pm1.15}$ & \textbf{56.95}$_{\pm0.55}$ \\
\bottomrule
\end{tabular}
}
\end{table}

\begin{table}[H]
\centering
\caption{Utility for \textbf{SPOUT} with \texttt{LLaMa-3.2-1B} embeddings. %\textit{Full Feedback} serves as an oracle upper bound. 
The best-performing method for each column is highlighted in bold.}
\label{tab:sprout_llama-3.2-1B}
\resizebox{\linewidth}{!}{%
\begin{tabular}{lccccccccccc}
\toprule
\textbf{Method} & $\lambda=0$ & $\lambda=100$ & $\lambda=200$ & $\lambda=300$ & $\lambda=400$ & $\lambda=500$ & $\lambda=600$ & $\lambda=700$ & $\lambda=800$ & $\lambda=900$ & $\lambda=1000$ \\
\midrule
Full-Feedback & \textit{85.19}$_{\pm0.17}$ & \textit{78.50}$_{\pm0.35}$ & \textit{74.47}$_{\pm0.28}$ & \textit{70.87}$_{\pm0.23}$ & \textit{67.44}$_{\pm0.22}$ & \textit{64.87}$_{\pm0.54}$ & 62.69$_{\pm0.36}$ & \textit{61.01}$_{\pm0.57}$ & \textit{59.51}$_{\pm0.55}$ & \textit{57.75}$_{\pm0.60}$ & \textit{56.22}$_{\pm0.58}$ \\ \hdashline
Baseline & 61.12$_{\pm1.95}$ & 57.12$_{\pm1.66}$ & 52.35$_{\pm2.63}$ & 51.23$_{\pm1.04}$ & 47.86$_{\pm2.19}$ & 42.50$_{\pm2.64}$ & 40.05$_{\pm4.79}$ & 34.20$_{\pm5.47}$ & 30.20$_{\pm5.82}$ & 29.64$_{\pm5.25}$ & 27.99$_{\pm5.66}$ \\
R\&C & 81.34$_{\pm0.41}$ & 74.75$_{\pm0.54}$ & 68.97$_{\pm5.81}$ & 67.09$_{\pm0.67}$ & 64.62$_{\pm0.57}$ & 62.33$_{\pm0.63}$ & 60.22$_{\pm0.46}$ & 58.81$_{\pm0.48}$ & 56.86$_{\pm0.81}$ & 55.76$_{\pm1.17}$ & 54.10$_{\pm1.48}$ \\
Causal-CARROT (kNN) & 84.50$_{\pm0.38}$ & 76.30$_{\pm0.42}$ & 71.24$_{\pm0.67}$ & 68.02$_{\pm0.54}$ & 65.57$_{\pm0.53}$ & 63.62$_{\pm0.47}$ & 61.93$_{\pm0.43}$ & 60.33$_{\pm0.41}$ & \textbf{59.05}$_{\pm0.37}$ & \textbf{57.86}$_{\pm0.35}$ & 56.89$_{\pm0.29}$ \\
Causal-CARROT (EmbedNet) & 81.29$_{\pm0.48}$ & 62.96$_{\pm1.32}$ & 55.69$_{\pm4.27}$ & 49.40$_{\pm5.41}$ & 45.21$_{\pm4.47}$ & 32.52$_{\pm14.00}$ & 31.05$_{\pm11.57}$ & 31.70$_{\pm6.75}$ & 22.34$_{\pm9.09}$ & 20.47$_{\pm9.64}$ & 15.52$_{\pm6.70}$ \\
CF-Regression & 82.32$_{\pm0.45}$ & 75.55$_{\pm0.43}$ & 70.24$_{\pm5.37}$ & 68.26$_{\pm0.64}$ & 65.55$_{\pm0.58}$ & 63.51$_{\pm0.52}$ & 61.49$_{\pm0.53}$ & 59.90$_{\pm0.60}$ & 58.19$_{\pm0.60}$ & 56.86$_{\pm0.61}$ & 55.61$_{\pm0.76}$ \\
RM-Classification & 83.60$_{\pm0.43}$ & 77.17$_{\pm0.47}$ & 70.95$_{\pm6.42}$ & 69.36$_{\pm0.62}$ & 66.38$_{\pm0.76}$ & 63.99$_{\pm0.89}$ & 61.82$_{\pm0.71}$ & 60.45$_{\pm0.68}$ & 58.38$_{\pm0.71}$ & 56.83$_{\pm0.83}$ & 55.64$_{\pm1.35}$ \\
RM-Softmax & \textbf{84.60}$_{\pm0.89}$ & \textbf{77.48}$_{\pm1.21}$ & \textbf{71.45}$_{\pm6.01}$ & 69.91$_{\pm0.85}$ & \textbf{67.18}$_{\pm0.51}$ & \textbf{65.35}$_{\pm0.60}$ & \textbf{63.04}$_{\pm1.35}$ & \textbf{61.38}$_{\pm1.61}$ & 59.04$_{\pm1.05}$ & 57.72$_{\pm0.33}$ & \textbf{57.12}$_{\pm0.21}$ \\ \hdashline
RM-Interval & \textbf{84.60}$_{\pm0.89}$ & 76.54$_{\pm4.06}$ & \textbf{71.45}$_{\pm6.01}$ & \textbf{69.99}$_{\pm0.62}$ & \textbf{67.18}$_{\pm0.51}$ & 64.73$_{\pm0.97}$ & \textbf{63.04}$_{\pm1.35}$ & 61.16$_{\pm1.84}$ & 59.04$_{\pm1.05}$ & 57.48$_{\pm2.90}$ & \textbf{57.12}$_{\pm0.21}$ \\
\bottomrule
\end{tabular}
}
\end{table}

We also report the AUC (Area Under the Curve) of the accuracy–cost trade-off curve for each method. This is computed using the \texttt{sklearn.metrics.auc}  function \cite{sklearn_api}. AUC provides a single scalar summary of performance that captures how well a model balances accuracy and computational cost. A higher AUC indicates a more favorable overall trade-off across budgets, making it a robust evaluation metric for comparing routing strategies.


\begin{table}[H]
\centering
\caption{Average AUC over 10 trials across datasets and embedding models. Higher is better. Abbreviations: RB = RouterBench, SP = SPROUT.}
\label{tab:single_lambda_auc}
\begin{tabular}{lcccc}
\toprule
\textbf{Method} & \textbf{RB-BERT} & \textbf{RB-LLaMa} & \textbf{SP-BERT} & \textbf{SP-LLaMa} \\
\midrule

Full-Feedback      & $0.1839$ & $0.1719$& $0.1727$  & $0.1655$  \\ \hdashline
Baseline    & $0.0890$ &  $0.0134$& $0.0255$ & $0.0148$ \\
R\&C      & $0.1539$ & $0.1108$ & $0.1316$ & $0.1105$ \\
Causal-CARROT (kNN)      & $0.1441$ & $0.1433$&$0.1642$  &  $0.1618$ \\
Causal-CARROT (EmbedNet)      & $0.1376$ & $0.0842$ & $0.1148$  & $0.0768$ \\


CF-Regression         & $0.1412$ & $0.1119$&  $0.1352$& $0.1233$ \\
RM-Classification     & $0.1665$ & $0.1389$ & $0.1477$ & $0.1436$ \\
RM-Softmax  & $\textbf{0.2286}$ & $\textbf{0.2213}$ &  $\textbf{0.3320}$ & $\textbf{0.2444}$  \\
\hdashline
RM-Interval  & $0.2285$ & $0.2196$ &  $0.3320$ & $0.2464$  \\
\bottomrule
\end{tabular}
\end{table}


\section{Contribution of Each Component}

The results highlight the impact of causal bias correction, end-to-end training, and regret-focused objectives as follows:


    \paragraph{Causal Inference (Bias Correction):} Our Baseline corresponds to CARROT without any causal correction - that is, a decoupled predictor trained directly on observational data without accounting for treatment bias. It consistently performs the worst across all settings. In contrast, incorporating causal techniques for bias correction yields substantial gains: both R\&C and Causal-CARROT, which integrate causal adjustments into routing, achieve +10–15\% routing accuracy at comparable cost, demonstrating that accounting for selection bias is critical for effective routing. These results validate that bias-aware routing significantly improves utility over naive predictors trained on biased data.
    \paragraph{End-to-End Learning vs. Two-Stage:} Among bias-corrected approaches, our end-to-end regret-minimizing methods (RM) consistently outperform the two-stage methods (CF-Regression, Causal-CARROT, R\&C). The performance gap (+1–3\% routing accuracy at comparable cost) demonstrates the benefit of an integrated end-to-end approach: by directly optimizing the decision-quality objective (regret) rather than optimizing intermediate predictions, our method achieves superior and more stable results.
    \paragraph{Regret Minimization Objective:} Even compared to other causal learners, our specific training objective provides an edge. For instance, RM-Softmax (differentiable surrogate) slightly outperforms RM-Classification (upper-bound surrogate) in most cases, with lower variance (+0.5–1\% routing accuracy at comparable cost). This highlights the advantage of our softmax-weighted regret objective and its alignment with the true decision loss. Moreover, both RM methods outperform Causal-CARROT and CF-Regression, underscoring that minimizing expected regret is more effective than surrogate approaches that focus only on accuracy/cost prediction.






\section{Experimental Details}\label{app:hyperparams}

All experiments were implemented in Python 3.8.12~\cite{python}, using PyTorch 2.4.1+cu121~\cite{NEURIPS2019_9015} and Scikit-learn~\cite{sklearn_api}. Experiments were conducted on an internal compute cluster equipped with an Intel(R) Xeon(R) Platinum 8260 CPU @ 2.40GHz, 512~GB of RAM, and two NVIDIA V100 GPUs with 16~GB memory each.

\paragraph{Prompt Encoding \& Augmentation} To encode input queries into vector representations \(x\), we employ a two-stage embedding process. First, we enrich each prompt with contextual metadata by prepending a natural language prefix that identifies the source dataset. Specifically, for a prompt \texttt{p} originating from dataset \texttt{D} (e.g., \texttt{openhermes/teknium}), we construct the following context-augmented input: \textit{``The following prompt comes from the dataset D. The prompt is: \texttt{p}''}. This step provides the embedding model with useful dataset-level context, which is particularly beneficial in multi-domain routing scenarios. The template is flexible and can be extended to include additional metadata if desired.

\paragraph{Datasets.} %We conduct experiments on two publicly available benchmark datasets for LLM routing: \textbf{RouterBench}~\citep{hu2024routerbench} and \textbf{SPROUT}~\citep{somerstep2025carrot}. 

\textbf{RouterBench}~\citep{hu2024routerbench} is a standardized benchmark for LLM routing, comprising 35,712 prompt-response pairs collected from 11 LLMs. The prompts span eight different evaluation benchmarks covering reasoning, factual knowledge, dialogue, mathematics, and code generation. Each prompt is annotated with model accuracy and execution cost, enabling response-based decision-making. To maintain consistency in evaluation, we adopt the same split strategy for RouterBench, applied deterministically at the prompt level to ensure reproducibility. \textbf{SPROUT}~\citep{somerstep2025carrot} is a more recent and larger benchmark for cost-aware routing, consisting of 44,241 prompts and responses from 13 state-of-the-art language models. The prompts are drawn from six diverse benchmarks, including GPQA \citep{rein2024gpqa}, MuSR \citep{sprague2023musr}, MMLU-Pro \citep{wang2024mmlu}, MATH \citep{hendrycks2021measuring}, OpenHermes \citep{teknium2023openhermes}, and RAGBench \citep{friel2407ragbench}. SPROUT includes a predefined train/validation/test split, using 80\% of the data for training and splitting the remaining 20\% equally between validation and test sets.


\paragraph{Neural Router Models.}  
All neural models used in our experiments share the same architecture for fairness and comparability. We use a 2-layer feedforward neural network with GELU activation and 200 hidden units per layer. Models are trained using the Adam optimizer with a learning rate of \(10^{-4}\), batch size of 128, and a maximum of 10,000 epochs. Early stopping is applied with a patience of 100 epochs based on validation regret. The temperature parameter for the softmax-based regret objective is set to 100, and to 1000 for the interval model to allow smoother gradients across budget intervals.

\paragraph{Doubly Robust Estimation.}  
For the outcome model \(\hat{r}_t(x)\), we use the same neural architecture described above, trained separately for each treatment \(t\). For the propensity model \(\hat{p}(t|x)\), we use an XGBoost classifier with the following hyperparameters: maximum depth = [1,2,3,5], number of trees = [10,20,50,100]. The estimated DR scores are clipped to the [5\textsuperscript{th}, 95\textsuperscript{th}] percentile to reduce the impact of extreme propensity weights and improve training stability.

\paragraph{Embedding Generation.}  
We generate sentence-level embeddings using the \texttt{bert-base-uncased} (768-dim) and \texttt{meta-llama/Llama-3.2-1B} (2048-dim) models. Embeddings are extracted via mean pooling over the final hidden states and are precomputed in batches using GPU acceleration. These embeddings are fixed during training of all downstream routing models.

\paragraph{RM-Interval Network.}  
The joint model used for budget interpolation is implemented using a small feedforward network that takes as input the concatenation of outputs from \(f_{\lambda_j}\), \(f_{\lambda_{j+1}}\), and a linear embedding of \(\lambda\). The architecture mirrors the router described above and is fine-tuned using the regret objective over interval-specific training data. The proposed architecture is presented in Figure \ref{fig:budget-aware-architecture}.

\paragraph{Inference latency and computational efficiency} Latency is critical in real-time applications. Our method is designed to minimize this by using lightweight routing networks (2-layer MLP with 200 neurons per hidden layer) and precomputed light-weight embeddings (Llama-3.2-1B and BERT). For instance, with Llama-3.2-1B embeddings and an MLP-based router on the RoutherBench dataset, the end-to-end routing latency is under 2.5 ms on a single A100 GPU for a batch size of almost 25,000. To contextualize this, recent benchmarks show that on 8× A100s, LLaMA‑2 or LLaMA‑3 70B can generate a ~100-token paragraph in 1.5 to 2.5 seconds under realistic conditions using optimized inference stacks. Even when using fewer GPUs or smaller models, generation latency typically remains one to two orders of magnitude higher than our routing time. In many cases (e.g., longer outputs or cold starts), the difference can be significantly larger. Thus, the routing overhead is negligible relative to the cost of LLM generation and does not meaningfully impact user experience. Additionally, the interval-conditioned architecture reduces deployment complexity by avoiding the need to train or store separate models for different user preferences.

% \newpage
% \section*{NeurIPS Paper Checklist}

% \begin{enumerate}

% \item {\bf Claims}
%     \item[] Question: Do the main claims made in the abstract and introduction accurately reflect the paper's contributions and scope?
%     \item[] Answer: \answerYes{}
%     \item[] Justification: The abstract and introduction accurately reflect the paper's contributions and scope. They clearly articulate the core innovations: (i) learning LLM routing policies from observational data, (ii) introducing an end-to-end regret minimization framework that avoids the limitations of decoupled approaches, (iii) developing differentiable surrogate objectives for training, and (iv) extending the method to support heterogeneous cost preferences via an interval-conditioned model. These contributions are substantiated in the main text, both theoretically and empirically, and are consistent with the stated goals of enabling scalable, bias-aware LLM routing in realistic deployment settings.

%     \item[] Guidelines:
%     \begin{itemize}
%         \item The answer NA means that the abstract and introduction do not include the claims made in the paper.
%         \item The abstract and/or introduction should clearly state the claims made, including the contributions made in the paper and important assumptions and limitations. A No or NA answer to this question will not be perceived well by the reviewers. 
%         \item The claims made should match theoretical and experimental results, and reflect how much the results can be expected to generalize to other settings. 
%         \item It is fine to include aspirational goals as motivation as long as it is clear that these goals are not attained by the paper. 
%     \end{itemize}

% \item {\bf Limitations}
%     \item[] Question: Does the paper discuss the limitations of the work performed by the authors?
%     \item[] Answer: \answerYes{} % Replace by \answerYes{}, \answerNo{}, or \answerNA{}.
%     \item[] Justification: The paper discusses its limitations explicitly. We note that the current framework is restricted to affine utility functions, which limits flexibility in modeling more complex user preferences (e.g., hard constraints that cannot be regularized in the objective). In the conclusion, we also outline directions for future work, including extending the method to dynamic or online learning settings and adapting it to multi-turn interactions. Additionally, the paper makes certain assumptions to support theoretical guarantees, including a mild Lipschitz continuity condition, and others that are standard in the causal inference literature (e.g., unconfoundedness). These are discussed in the relevant sections and are typical for counterfactual utility estimation from observational data.

%     \item[] Guidelines:
%     \begin{itemize}
%         \item The answer NA means that the paper has no limitation while the answer No means that the paper has limitations, but those are not discussed in the paper. 
%         \item The authors are encouraged to create a separate "Limitations" section in their paper.
%         \item The paper should point out any strong assumptions and how robust the results are to violations of these assumptions (e.g., independence assumptions, noiseless settings, model well-specification, asymptotic approximations only holding locally). The authors should reflect on how these assumptions might be violated in practice and what the implications would be.
%         \item The authors should reflect on the scope of the claims made, e.g., if the approach was only tested on a few datasets or with a few runs. In general, empirical results often depend on implicit assumptions, which should be articulated.
%         \item The authors should reflect on the factors that influence the performance of the approach. For example, a facial recognition algorithm may perform poorly when image resolution is low or images are taken in low lighting. Or a speech-to-text system might not be used reliably to provide closed captions for online lectures because it fails to handle technical jargon.
%         \item The authors should discuss the computational efficiency of the proposed algorithms and how they scale with dataset size.
%         \item If applicable, the authors should discuss possible limitations of their approach to address problems of privacy and fairness.
%         \item While the authors might fear that complete honesty about limitations might be used by reviewers as grounds for rejection, a worse outcome might be that reviewers discover limitations that aren't acknowledged in the paper. The authors should use their best judgment and recognize that individual actions in favor of transparency play an important role in developing norms that preserve the integrity of the community. Reviewers will be specifically instructed to not penalize honesty concerning limitations.
%     \end{itemize}

% \item {\bf Theory assumptions and proofs}
%     \item[] Question: For each theoretical result, does the paper provide the full set of assumptions and a complete (and correct) proof?
%     \item[] Answer: \answerYes{} %\answerTODO{} % Replace by \answerYes{}, \answerNo{}, or \answerNA{}.
%     \item[] Justification: All theorems, assumptions, propositions, and corollaries are included in the paper. Their complete proofs can be found in the Appendix.
%     \item[] Guidelines:
%     \begin{itemize}
%         \item The answer NA means that the paper does not include theoretical results. 
%         \item All the theorems, formulas, and proofs in the paper should be numbered and cross-referenced.
%         \item All assumptions should be clearly stated or referenced in the statement of any theorems.
%         \item The proofs can either appear in the main paper or the supplemental material, but if they appear in the supplemental material, the authors are encouraged to provide a short proof sketch to provide intuition. 
%         \item Inversely, any informal proof provided in the core of the paper should be complemented by formal proofs provided in appendix or supplemental material.
%         \item Theorems and Lemmas that the proof relies upon should be properly referenced. 
%     \end{itemize}

%     \item {\bf Experimental result reproducibility}
%     \item[] Question: Does the paper fully disclose all the information needed to reproduce the main experimental results of the paper to the extent that it affects the main claims and/or conclusions of the paper (regardless of whether the code and data are provided or not)?
%     \item[] Answer: \answerYes{}
%     \item[] Justification: You can find all the details to reproduce the experiments in Section \ref{sec:exp} and in Section \ref{app:hyperparams} of the Appendix. We will also release code in the supplementary material. 
%     \item[] Guidelines:
%     \begin{itemize}
%         \item The answer NA means that the paper does not include experiments.
%         \item If the paper includes experiments, a No answer to this question will not be perceived well by the reviewers: Making the paper reproducible is important, regardless of whether the code and data are provided or not.
%         \item If the contribution is a dataset and/or model, the authors should describe the steps taken to make their results reproducible or verifiable. 
%         \item Depending on the contribution, reproducibility can be accomplished in various ways. For example, if the contribution is a novel architecture, describing the architecture fully might suffice, or if the contribution is a specific model and empirical evaluation, it may be necessary to either make it possible for others to replicate the model with the same dataset, or provide access to the model. In general. releasing code and data is often one good way to accomplish this, but reproducibility can also be provided via detailed instructions for how to replicate the results, access to a hosted model (e.g., in the case of a large language model), releasing of a model checkpoint, or other means that are appropriate to the research performed.
%         \item While NeurIPS does not require releasing code, the conference does require all submissions to provide some reasonable avenue for reproducibility, which may depend on the nature of the contribution. For example
%         \begin{enumerate}
%             \item If the contribution is primarily a new algorithm, the paper should make it clear how to reproduce that algorithm.
%             \item If the contribution is primarily a new model architecture, the paper should describe the architecture clearly and fully.
%             \item If the contribution is a new model (e.g., a large language model), then there should either be a way to access this model for reproducing the results or a way to reproduce the model (e.g., with an open-source dataset or instructions for how to construct the dataset).
%             \item We recognize that reproducibility may be tricky in some cases, in which case authors are welcome to describe the particular way they provide for reproducibility. In the case of closed-source models, it may be that access to the model is limited in some way (e.g., to registered users), but it should be possible for other researchers to have some path to reproducing or verifying the results.
%         \end{enumerate}
%     \end{itemize}


% \item {\bf Open access to data and code}
%     \item[] Question: Does the paper provide open access to the data and code, with sufficient instructions to faithfully reproduce the main experimental results, as described in supplemental material?
%     \item[] Answer: \answerYes{} % Replace by \answerYes{}, \answerNo{}, or \answerNA{}.
%     \item[] Justification: Our paper provides all the details to reproduce the experimental results in Section \ref{sec:exp} and in Section \ref{app:hyperparams} of the Appendix. Furthermore, the code used is provided in the supplemental materials.
%     \item[] Guidelines:
%     \begin{itemize}
%         \item The answer NA means that paper does not include experiments requiring code.
%         \item Please see the NeurIPS code and data submission guidelines (\url{https://nips.cc/public/guides/CodeSubmissionPolicy}) for more details.
%         \item While we encourage the release of code and data, we understand that this might not be possible, so “No” is an acceptable answer. Papers cannot be rejected simply for not including code, unless this is central to the contribution (e.g., for a new open-source benchmark).
%         \item The instructions should contain the exact command and environment needed to run to reproduce the results. See the NeurIPS code and data submission guidelines (\url{https://nips.cc/public/guides/CodeSubmissionPolicy}) for more details.
%         \item The authors should provide instructions on data access and preparation, including how to access the raw data, preprocessed data, intermediate data, and generated data, etc.
%         \item The authors should provide scripts to reproduce all experimental results for the new proposed method and baselines. If only a subset of experiments are reproducible, they should state which ones are omitted from the script and why.
%         \item At submission time, to preserve anonymity, the authors should release anonymized versions (if applicable).
%         \item Providing as much information as possible in supplemental material (appended to the paper) is recommended, but including URLs to data and code is permitted.
%     \end{itemize}


% \item {\bf Experimental setting/details}
%     \item[] Question: Does the paper specify all the training and test details (e.g., data splits, hyperparameters, how they were chosen, type of optimizer, etc.) necessary to understand the results?
%     \item[] Answer: \answerYes{} % Replace by \answerYes{}, \answerNo{}, or \answerNA{}.
%     \item[] Justification: We report all training and evaluation details in Section~\ref{sec:exp} and Appendix~\ref{app:hyperparams}.
%     \item[] Guidelines:
%     \begin{itemize}
%         \item The answer NA means that the paper does not include experiments.
%         \item The experimental setting should be presented in the core of the paper to a level of detail that is necessary to appreciate the results and make sense of them.
%         \item The full details can be provided either with the code, in appendix, or as supplemental material.
%     \end{itemize}

% \item {\bf Experiment statistical significance}
%     \item[] Question: Does the paper report error bars suitably and correctly defined or other appropriate information about the statistical significance of the experiments?
%     \item[] Answer: \answerYes{} % Replace by \answerYes{}, \answerNo{}, or \answerNA{}.
%     \item[] Justification: We run each experiment 10 times and report mean utility along with standard deviation (1-sigma). The variability reflects random initialization and train/validation splits. 
%     \item[] Guidelines:
%     \begin{itemize}
%         \item The answer NA means that the paper does not include experiments.
%         \item The authors should answer "Yes" if the results are accompanied by error bars, confidence intervals, or statistical significance tests, at least for the experiments that support the main claims of the paper.
%         \item The factors of variability that the error bars are capturing should be clearly stated (for example, train/test split, initialization, random drawing of some parameter, or overall run with given experimental conditions).
%         \item The method for calculating the error bars should be explained (closed form formula, call to a library function, bootstrap, etc.)
%         \item The assumptions made should be given (e.g., Normally distributed errors).
%         \item It should be clear whether the error bar is the standard deviation or the standard error of the mean.
%         \item It is OK to report 1-sigma error bars, but one should state it. The authors should preferably report a 2-sigma error bar than state that they have a 96\% CI, if the hypothesis of Normality of errors is not verified.
%         \item For asymmetric distributions, the authors should be careful not to show in tables or figures symmetric error bars that would yield results that are out of range (e.g. negative error rates).
%         \item If error bars are reported in tables or plots, The authors should explain in the text how they were calculated and reference the corresponding figures or tables in the text.
%     \end{itemize}

% \item {\bf Experiments compute resources}
%     \item[] Question: For each experiment, does the paper provide sufficient information on the computer resources (type of compute workers, memory, time of execution) needed to reproduce the experiments?
%     \item[] Answer: \answerYes{}
%     \item[] Justification: We state the computer resources used in Section \ref{app:hyperparams} of the Appendix.
%     \item[] Guidelines:
%     \begin{itemize}
%         \item The answer NA means that the paper does not include experiments.
%         \item The paper should indicate the type of compute workers CPU or GPU, internal cluster, or cloud provider, including relevant memory and storage.
%         \item The paper should provide the amount of compute required for each of the individual experimental runs as well as estimate the total compute. 
%         \item The paper should disclose whether the full research project required more compute than the experiments reported in the paper (e.g., preliminary or failed experiments that didn't make it into the paper). 
%     \end{itemize}
    
% \item {\bf Code of ethics}
%     \item[] Question: Does the research conducted in the paper conform, in every respect, with the NeurIPS Code of Ethics \url{https://neurips.cc/public/EthicsGuidelines}?
%     \item[] Answer: \answerYes{} % Replace by \answerYes{}, \answerNo{}, or \answerNA{}.
%     \item[] Justification: We have reviewed the NeurIPS Code of Ethics, and our work complies with its guidelines in all respects.
%     \item[] Guidelines:
%     \begin{itemize}
%         \item The answer NA means that the authors have not reviewed the NeurIPS Code of Ethics.
%         \item If the authors answer No, they should explain the special circumstances that require a deviation from the Code of Ethics.
%         \item The authors should make sure to preserve anonymity (e.g., if there is a special consideration due to laws or regulations in their jurisdiction).
%     \end{itemize}


% \item {\bf Broader impacts}
%     \item[] Question: Does the paper discuss both potential positive societal impacts and negative societal impacts of the work performed?
%     \item[] Answer: \answerYes{}% Replace by \answerYes{}, \answerNo{}, or \answerNA{}.
%     \item[] Justification: The paper highlights potential positive societal impacts of this work. By enabling efficient LLM routing using observational data, our framework reduces the need for full-feedback supervision, significantly lowering computational and monetary costs. LLM routing itself contributes to more sustainable AI deployment by reducing energy consumption and resource usage. We do not anticipate negative societal impacts, as the method focuses on infrastructure-level efficiency improvements and does not involve sensitive user data or content generation.

%     \item[] Guidelines:
%     \begin{itemize}
%         \item The answer NA means that there is no societal impact of the work performed.
%         \item If the authors answer NA or No, they should explain why their work has no societal impact or why the paper does not address societal impact.
%         \item Examples of negative societal impacts include potential malicious or unintended uses (e.g., disinformation, generating fake profiles, surveillance), fairness considerations (e.g., deployment of technologies that could make decisions that unfairly impact specific groups), privacy considerations, and security considerations.
%         \item The conference expects that many papers will be foundational research and not tied to particular applications, let alone deployments. However, if there is a direct path to any negative applications, the authors should point it out. For example, it is legitimate to point out that an improvement in the quality of generative models could be used to generate deepfakes for disinformation. On the other hand, it is not needed to point out that a generic algorithm for optimizing neural networks could enable people to train models that generate Deepfakes faster.
%         \item The authors should consider possible harms that could arise when the technology is being used as intended and functioning correctly, harms that could arise when the technology is being used as intended but gives incorrect results, and harms following from (intentional or unintentional) misuse of the technology.
%         \item If there are negative societal impacts, the authors could also discuss possible mitigation strategies (e.g., gated release of models, providing defenses in addition to attacks, mechanisms for monitoring misuse, mechanisms to monitor how a system learns from feedback over time, improving the efficiency and accessibility of ML).
%     \end{itemize}
    
% \item {\bf Safeguards}
%     \item[] Question: Does the paper describe safeguards that have been put in place for responsible release of data or models that have a high risk for misuse (e.g., pretrained language models, image generators, or scraped datasets)?
%     \item[] Answer: \answerYes{} %, \answerNo{}, or \answerNA{}.
%     \item[] Justification: Yes. The paper uses publicly available benchmark datasets, and we plan to release our code to support transparency and reproducibility. As our work focuses on routing policies rather than generating model outputs, and does not involve training or releasing new pretrained language models or sensitive datasets, we do not identify any specific risks or misuse concerns requiring additional safeguards.

%     \item[] Guidelines:
%     \begin{itemize}
%         \item The answer NA means that the paper poses no such risks.
%         \item Released models that have a high risk for misuse or dual-use should be released with necessary safeguards to allow for controlled use of the model, for example by requiring that users adhere to usage guidelines or restrictions to access the model or implementing safety filters. 
%         \item Datasets that have been scraped from the Internet could pose safety risks. The authors should describe how they avoided releasing unsafe images.
%         \item We recognize that providing effective safeguards is challenging, and many papers do not require this, but we encourage authors to take this into account and make a best faith effort.
%     \end{itemize}

% \item {\bf Licenses for existing assets}
%     \item[] Question: Are the creators or original owners of assets (e.g., code, data, models), used in the paper, properly credited and are the license and terms of use explicitly mentioned and properly respected?
%     \item[] Answer: \answerYes{} %, \answerNo{}, or \answerNA{}.
%     \item[] Justification: All external assets used in the paper, including benchmark datasets, embedding models, are properly cited in the main text. We have respected the terms of use and licenses associated with these resources, and references are provided to acknowledge the original creators.

%     \item[] Guidelines:
%     \begin{itemize}
%         \item The answer NA means that the paper does not use existing assets.
%         \item The authors should cite the original paper that produced the code package or dataset.
%         \item The authors should state which version of the asset is used and, if possible, include a URL.
%         \item The name of the license (e.g., CC-BY 4.0) should be included for each asset.
%         \item For scraped data from a particular source (e.g., website), the copyright and terms of service of that source should be provided.
%         \item If assets are released, the license, copyright information, and terms of use in the package should be provided. For popular datasets, \url{paperswithcode.com/datasets} has curated licenses for some datasets. Their licensing guide can help determine the license of a dataset.
%         \item For existing datasets that are re-packaged, both the original license and the license of the derived asset (if it has changed) should be provided.
%         \item If this information is not available online, the authors are encouraged to reach out to the asset's creators.
%     \end{itemize}

% \item {\bf New assets}
%     \item[] Question: Are new assets introduced in the paper well documented and is the documentation provided alongside the assets?
%     \item[] Answer: \answerYes{} %, \answerNo{}, or \answerNA{}.
%     \item[] Justification: We introduce new code assets to implement our causal end-to-end LLM routing framework. These assets are well documented, with usage instructions and reproducibility details provided to facilitate adoption and verification. Documentation will be released alongside the code repository.

%     \item[] Guidelines:
%     \begin{itemize}
%         \item The answer NA means that the paper does not release new assets.
%         \item Researchers should communicate the details of the dataset/code/model as part of their submissions via structured templates. This includes details about training, license, limitations, etc. 
%         \item The paper should discuss whether and how consent was obtained from people whose asset is used.
%         \item At submission time, remember to anonymize your assets (if applicable). You can either create an anonymized URL or include an anonymized zip file.
%     \end{itemize}

% \item {\bf Crowdsourcing and research with human subjects}
%     \item[] Question: For crowdsourcing experiments and research with human subjects, does the paper include the full text of instructions given to participants and screenshots, if applicable, as well as details about compensation (if any)? 
%     \item[] Answer:  \answerNA{}
%     \item[] Justification: We did not conduct crowdsourcing experiments and research with human subjects.
%     \item[] Guidelines:
%     \begin{itemize}
%         \item The answer NA means that the paper does not involve crowdsourcing nor research with human subjects.
%         \item Including this information in the supplemental material is fine, but if the main contribution of the paper involves human subjects, then as much detail as possible should be included in the main paper. 
%         \item According to the NeurIPS Code of Ethics, workers involved in data collection, curation, or other labor should be paid at least the minimum wage in the country of the data collector. 
%     \end{itemize}

% \item {\bf Institutional review board (IRB) approvals or equivalent for research with human subjects}
%     \item[] Question: Does the paper describe potential risks incurred by study participants, whether such risks were disclosed to the subjects, and whether Institutional Review Board (IRB) approvals (or an equivalent approval/review based on the requirements of your country or institution) were obtained?
%     \item[] Answer: \answerNA{}
%     \item[] Justification: The paper does not involve crowdsourcing nor research with human subjects.
%     \item[] Guidelines:
%     \begin{itemize}
%         \item The answer NA means that the paper does not involve crowdsourcing nor research with human subjects.
%         \item Depending on the country in which research is conducted, IRB approval (or equivalent) may be required for any human subjects research. If you obtained IRB approval, you should clearly state this in the paper. 
%         \item We recognize that the procedures for this may vary significantly between institutions and locations, and we expect authors to adhere to the NeurIPS Code of Ethics and the guidelines for their institution. 
%         \item For initial submissions, do not include any information that would break anonymity (if applicable), such as the institution conducting the review.
%     \end{itemize}

% \item {\bf Declaration of LLM usage}
%     \item[] Question: Does the paper describe the usage of LLMs if it is an important, original, or non-standard component of the core methods in this research? Note that if the LLM is used only for writing, editing, or formatting purposes and does not impact the core methodology, scientific rigorousness, or originality of the research, declaration is not required.
%     %this research? 
%     \item[] Answer: \answerNA{}  % Replace by \answerYes{}, \answerNo{}, or \answerNA{}.
%     \item[] Justification: We used LLMs for grammar, formatting, and rephrasing assistance (e.g., improving proof presentation). They did not influence the methodological originality of the work.
%     \item[] Guidelines:
%     \begin{itemize}
%         \item The answer NA means that the core method development in this research does not involve LLMs as any important, original, or non-standard components.
%         \item Please refer to our LLM policy (\url{https://neurips.cc/Conferences/2025/LLM}) for what should or should not be described.
%     \end{itemize}

% \end{enumerate}

\end{document}
